% !TeX root = ../../Book2.tex
\chapter{多元分析与偏微分方程预备}
\label{b2:ch:01}
\BookChapterNumber{b2:ch:01}
% 本章摘要（不计入编号）
\begin{chapterabstract}
多元微积分中的局部计算与积分公式，是本章研究偏微分方程的出发点。先从方向导数和全微分出发，建立链式法则、隐函数描述及常用微分算子；再固定换元、面积元和边界定向，整理 Gauss、Green 与 Stokes 公式。含参变量积分一节说明连续性、求导和无穷尾部各需要什么控制。最后通过输运方程和 Poisson 方程介绍初值、边值及不同意义的解。后续 Fourier 分析和方程求解将用到这些概念与公式。
\end{chapterabstract}

% 本章目标（不计入编号）
\begin{learninggoals}
\begin{itemize}
\item\label{b2:item:01-goal-differential}能区分方向导数、Fréchet 微分与连续可微，利用链式法则和隐函数定理计算，并辨认局部结论的范围。
\item\label{b2:item:01-goal-vector}能为换元和曲面积分选取合法参数，核对法向与边界方向，使用 Green 公式转移导数。
\item\label{b2:item:01-goal-parameter}能为参数极限和积分号下求导找到局部统一控制，区分绝对可积与靠振荡收敛的反常积分。
\item\label{b2:item:01-goal-pde}能判断方程的阶数和非线性类型，写清初值或边值及解的意义，并用能量恒等式检验唯一性与相容条件。
\end{itemize}
\end{learninggoals}

考虑一个区域中稳态温度的分布。先取常数热导率 \(\kappa>0\) 的均匀介质，热通量向量为 \(q=-\kappa\nabla u\)：热量沿温度降低的方向流动。若单位体积的热源强度为 \(s\)，稳态平衡给出 \(\operatorname{div}q=s\)，即 \(-\kappa\Delta u=s\)。下文采用归一化变量，将 \(q/\kappa\)、\(s/\kappa\) 分别仍记为 \(q\)、\(f\)，于是 \(q=-\nabla u\)，无源时为 \(\Delta u=0\)，有源时为 \(-\Delta u=f\)。在此模型中，外向热通量是 \(q\cdot\nu=-\partial_\nu u\)，负号不能省去。

\ProseParagraphBreak

指定边界温度与指定流出边界的热通量，是两种不同的问题。要从方程读出总热量的平衡，须把内部的散度积分变成边界通量；要比较两个温度分布，须把二阶导数转移为非负的梯度能量。如果介质在界面两侧不同，温度的斜率还可能发生跳跃，逐点二阶导数便不再是合适的描述。多元微分、向量积分与弱解在这个问题中各自承担了明确的工作。

\ProseParagraphBreak

本章预设一元微积分、有限维线性代数和 Lebesgue 积分基础。第一卷中已经完整证明的逆函数、换元、面积与收敛定理，在本章首次引用时写明卷号及编号；读者可据此回查证明。弱导数与 Sobolev 空间在\subsecref{b2:subsec:010405}回顾定义和本章所需的解释，其一般理论仍以第一卷为基础。

\ProseParagraphBreak

\secref{b2:sec:0101}与\secref{b2:sec:0102}介绍多元微分与积分公式，\secref{b2:sec:0103}讨论含参积分的极限与求导，\secref{b2:sec:0104}介绍初值、边值以及不同意义的解。熟悉多元微积分的读者可以先核对本章的符号和边界约定，再集中阅读含参变量积分的尾部判据及强解、弱解的区别。首次学习时宜先完成普通习题，再回到带星号的边界能量题。

\ProseParagraphBreak

Sheldon Axler 的《Measure, Integration \& Real Analysis》\cite[Chapter 3]{Axler2020Measure}讨论积分与极限，阅读时可比较非负逼近和可积控制的不同作用。John K. Hunter 与 Bruno Nachtergaele 的《Applied Analysis》\cite[Section 12.9]{Hunter2001Applied}介绍弱导数，可将其中的一维积分恒等式与这里的分部积分表达相互对照。

% 本章来源实读范围见参考文献原文件/来源记录.md；证明义务及后续复核以现行正文和工作记录/审核进度.md为准。
% 第一卷接口依当前源文件及 Book1.aux 核对，逐项登记于跨卷引用表.md。
% Hunter2001Applied Section12.9，作者公开ch12.pdf，印刷362--364页；
% Axler2020Measure Chapter3 积分及收敛定理。目录核名与数学实读范围分开记录。

% !TeX root = ../../Book2.tex
\section{多元微分学}
\label{b2:sec:0101}

% 数学依赖：第一卷定义 16.2.2，b1:def:frechet-differentiability；
% 第一卷引理 18.1.2，b1:lem:level-set-local-inverse，完整压缩迭代证明。
% 本节其余结果按一元微积分与有限维线性代数独立展开，反例与计算自行组织。

偏微分方程把一个函数在不同方向上的变化联系起来。要正确读取这样的关系，需要先分清三件事：沿固定直线能否求导，小增量能否由同一个线性映射逼近，以及导数在邻近点之间是否连续。本节在开集上讨论经典微分；边界上的导数与较弱意义的导数将在需要时另行说明。以下向量使用欧氏范数 \(\lvert\cdot\rvert\)，矩阵使用由它诱导的算子范数 \(\lVert\cdot\rVert\)。

\subsection{方向导数}
\label{b2:subsec:010101}

\begin{definition}\label{b2:def:01-directional}
设 \(U\subseteq\mathbb R^d\) 为开集，\(f:U\rightarrow\mathbb R^m\)，且 \(a\in U\)。给定 \(v\in\mathbb R^d\)，若有限极限
\[
 D_vf(a)\coloneqq\lim_{t\to0}\frac{f(a+tv)-f(a)}{t}
\]
存在，则称它为 \(f\) 在 \(a\) 处沿 \(v\) 的\term[fangxiang-daoshu]{方向导数}[directional derivative]。取第 \(j\) 个坐标向量 \(e_j\)，所得 \(D_{e_j}f(a)\) 称为第 \(j\) 个\term[piandaoshu]{偏导数}[partial derivative]，记作 \(\partial_jf(a)\)。
\end{definition}
\symidx{\(D_vf\)：沿向量 \(v\) 的方向导数}
\symidx{\(\partial_jf\)：第 \(j\) 个偏导数}

这里不要求 \(v\) 为单位向量；当 \(\lvert v\rvert=1\) 时，参数 \(t\) 才直接计量沿直线的有向距离。定义中的极限为双侧极限，开集条件保证充分小的正、负 \(t\) 都有意义。例如 \(u(x,y)=x^2+xy\) 满足
\[
 D_{(a,b)}u(x,y)=a(2x+y)+bx.
\]
右端关于方向 \((a,b)\) 线性。不过，这一性质不能从方向导数的存在本身推出。

\ProseParagraphBreak

事实上，即使所有方向导数都为零，也不能断言函数连续。令
\[
 f(x,y)=
 \begin{cases}
 \dfrac{x^3y}{x^6+y^2},&(x,y)\ne(0,0),\\[4pt]
 0,&(x,y)=(0,0).
 \end{cases}
\]
若 \(b\ne0\)，沿方向 \((a,b)\) 有
\[
 \frac{f(ta,tb)}t=\frac{ta^3b}{t^4a^6+b^2}\longrightarrow0;
\]
若 \(b=0\)，差商恒为零。然而 \(f(x,x^3)=1/2\) 对一切 \(x\ne0\) 成立。直线逐条检查的极限，没有控制沿弯曲路径趋近的行为。研究坐标变换和微分方程时，需要一个同时约束全部小增量的定义。

\subsection{Fréchet 微分}
\label{b2:subsec:010102}

\begin{definition}\label{b2:def:01-frechet}
设 \(U\subseteq\mathbb R^d\) 为开集，\(f:U\rightarrow\mathbb R^m\)，且 \(a\in U\)。若存在线性映射 \(A:\mathbb R^d\rightarrow\mathbb R^m\)，使
\[
 \lim_{\substack{h\to0\\h\ne0}}
 \frac{\left|f(a+h)-f(a)-Ah\right|}{\lvert h\rvert}=0,
\]
则称 \(f\) 在 \(a\) 处\term[frechet-kewei]{Fréchet 可微}[Fréchet differentiable]，称 \(A\) 为其\term[frechet-weifen]{Fréchet 微分}[Fréchet differential]，记作 \(Df(a)\)。在有限维空间中，这就是\term[quan-weifen]{全微分}[total differential]；下文未加限定的“可微”均指这一意义。
\end{definition}
\symidx{\(Df(a)\)：映射在一点的 Fréchet 微分}

这与第一卷定义 16.2.2 的全微分约定一致。余项也可记为 \(o(\lvert h\rvert)\)，意思是其范数除以 \(\lvert h\rvert\) 趋于零。由于线性映射在有限维空间中连续，可微性立即给出 \(f(a+h)\to f(a)\)。令 \(h=tv\) 又得 \(D_vf(a)=Df(a)v\)，所以微分由各坐标偏导唯一确定。设 \(U\subset\mathbb R^d\) 开，\(a\in U\)，\(f:U\to\mathbb R^m\) 在 \(a\) 处可微。若 \(f=(f_1,\ldots,f_m)\)，其矩阵表示
\[
 Df(a)=\left(\partial_jf_i(a)\right)_{1\leq i\leq m,\,1\leq j\leq d}
\]
称为\term[jacobi-juzhen]{Jacobi 矩阵}[Jacobian matrix]。这里专指由一阶偏导数组成的导数矩阵，与某些代数文献中同名的三对角矩阵不同。行指标对应输出分量，列指标对应输入坐标；方阵时的 \(\det Df(a)\) 称为\term[jacobi-hanglieshi]{Jacobi 行列式}[Jacobian determinant]，也称雅可比行列式。\termidx[yakebi-hanglieshi]{雅可比行列式}\footnote{萧树铁、扈志明《微积分（下）修订版》的\href{https://www.tup.com.cn/booksCenter/book_02932206.html}{出版社目录}第9.3.1节以“变量代换的雅可比行列式”为题。本书主称保留人名原拼，使用 Jacobi 行列式。}

\ProseParagraphBreak

实际计算中，偏导数的邻域连续性提供了一个便于核验的充分条件。设 \(U\subset\mathbb R^d\) 为开集。记 \(C^1(U;\mathbb R^m)\) 为一阶偏导数存在且连续的映射类；这样的映射称为\term[lianxu-kewei]{连续可微}[continuously differentiable]。记号 \(C^k\) 表示直到 \(k\) 阶的各阶偏导数都连续，\(C^\infty\) 表示任意有限阶偏导数都连续；这些记号本身不要求全域有界。
\symidx{\(C^k(U;\mathbb R^m)\)：具有连续偏导数直到 \(k\) 阶的映射类}

\begin{proposition}\label{b2:prop:01-c1-criterion}
设 \(U\subseteq\mathbb R^d\) 为开集，\(f:U\to\mathbb R^m\)、\(a\in U\)。若 \(f\) 的各一阶偏导数在 \(a\) 的某邻域内存在，且在 \(a\) 连续，则 \(f\) 在 \(a\) 可微。特别地，\(C^1\) 映射处处可微，且 \(x\mapsto Df(x)\) 连续。
\end{proposition}
\begin{proof}
先设 \(f\) 实值。把从 \(a\) 到 \(a+h\) 的位移拆成依次改变坐标的 \(d\) 段。每一段上的一元函数可微，中值定理给出该段增量为 \(h_j\cdot\partial_jf(\xi_j)\)，其中 \(\lvert\xi_j-a\rvert\leq\lvert h\rvert\)；当 \(h_j=0\) 时该项直接为零。因此
\[
 \left|f(a+h)-f(a)-\sum_{j=1}^d h_j\cdot\partial_jf(a)\right|
 \leq\sum_{j=1}^d\lvert h_j\rvert\cdot
 \left|\partial_jf(\xi_j)-\partial_jf(a)\right|
 =o(\lvert h\rvert),
\]
其中用了 \(\sum_j\lvert h_j\rvert\leq\sqrt d\,\lvert h\rvert\)。向量值情形对有限个分量分别应用此结论，余项的欧氏范数仍为 \(o(\lvert h\rvert)\)。最后，矩阵的有限个系数连续就给出算子范数意义的连续性。
\end{proof}

连续偏导是充分条件，不能误读为可微的必要条件。一元函数 \(u(t)=t^2\sin(1/t)\) 在 \(t\ne0\) 时按此式取值，另设 \(u(0)=0\)，则 \(u'(0)=0\)，但
\[
 u'(t)=2t\sin(1/t)-\cos(1/t)\quad(t\ne0)
\]
在零点没有极限。可微只比较函数与一点处的线性模型；连续可微还约束这些模型随基点如何变化。

\ProseParagraphBreak

\term[lianshi-faze]{链式法则}[chain rule]把线性近似与函数复合联系起来。它只要求涉及的两点可微，并不要求整个邻域都是 \(C^1\)。

\begin{proposition}[链式法则]\label{b2:prop:01-chain-rule}
设 \(f:U\rightarrow V\subseteq\mathbb R^m\)、\(g:V\rightarrow\mathbb R^n\)，其中 \(U\subseteq\mathbb R^d\)、\(V\subseteq\mathbb R^m\) 为开集，\(a\in U\)。若 \(f\) 在 \(a\) 可微，\(g\) 在 \(f(a)\) 可微，则
\[
 D(g\circ f)(a)=Dg\bigl(f(a)\bigr)\circ Df(a).
\]
\end{proposition}
\begin{proof}
令 \(A=Df(a)\)、\(B=Dg\bigl(f(a)\bigr)\)，写
\[
 f(a+h)-f(a)=Ah+r(h),\quad
 g\bigl(f(a)+k\bigr)-g\bigl(f(a)\bigr)=Bk+s(k),
\]
其中 \(r(h)=o(\lvert h\rvert)\)、\(s(k)=o(\lvert k\rvert)\)，并在零点把余项取为零。代入 \(k=Ah+r(h)\)，有 \(\lvert k\rvert\leq C\lvert h\rvert\)，复合的余项为 \(Br(h)+s(k)=o(\lvert h\rvert)\)。这里对 \(s(k)\) 的估计在 \(k=0\) 时也成立。
\end{proof}

一阶微分决定切向变化，二阶微分记录弯曲。设 \(U\subset\mathbb R^d\) 开、\(a\in U\)。对实值 \(u\in C^2(U)\)，矩阵
\[
 D^2u(a)\coloneqq\left(\partial_i\partial_ju(a)\right)_{i,j=1}^d
\]
称为\term[hesse-juzhen]{Hesse 矩阵}[Hessian matrix]。
\symidx{\(D^2u\)：实值函数的 Hesse 矩阵}
下面的\term[taylor-gongshi]{Taylor 公式}[Taylor formula]将其与函数增量联系起来。

\begin{proposition}[二阶 Taylor 公式]\label{b2:prop:01-taylor}
设 \(U\subseteq\mathbb R^d\) 为开集，\(u\in C^2(U;\mathbb R)\)、\(a\in U\)。则 \(D^2u(a)\) 为对称矩阵。对包含于 \(U\) 的线段 \([a,a+h]\)，有
\[
 u(a+h)=u(a)+Du(a)h+
 \int_0^1(1-t)h^{\mathsf T}D^2u(a+th)h\,\mathrm{d}t.
\]
因而在 \(h\to0\) 时，最后一项等于
\(\dfrac12 h^{\mathsf T}D^2u(a)h+o(\lvert h\rvert^2)\)。
\end{proposition}
\begin{proof}
固定 \(i\ne j\)，对充分小的正数 \(s,t\)，考虑矩形的交错增量
\[
 R(s,t)=u(a+se_i+te_j)-u(a+se_i)-u(a+te_j)+u(a).
\]
先沿 \(i\) 方向、再沿 \(j\) 方向用一元微积分基本定理，得到
\[
 R(s,t)=\int_0^s\int_0^t
 \partial_j\partial_i u(a+\sigma e_i+\tau e_j)
 \,\mathrm{d}\tau\,\mathrm{d}\sigma.
\]
改变两次求增量的次序，则得到含 \(\partial_i\partial_j u\) 的相应迭代积分。分别除以 \(st\)，由被积函数在 \(a\) 连续，两式在 \(s,t\downarrow0\) 时趋于相等的混合偏导数。

令 \(q(t)=u(a+th)\)。由链式法则，\(q'(0)=Du(a)h\)、\(q''(t)=h^{\mathsf T}D^2u(a+th)h\)。对 \(\int_0^1(1-t)q''(t)\,\mathrm{d}t\) 作一元分部积分即得公式。最后，积分余项减去 \(h^{\mathsf T}D^2u(a)h/2\) 的绝对值不超过
\[
 \frac{\lvert h\rvert^2}{2}
 \sup_{0\leq t\leq1}\left\lVert D^2u(a+th)-D^2u(a)\right\rVert,
\]
而上确界随 \(h\to0\) 趋于零。
\end{proof}

Taylor 公式还说明二阶项能在什么范围内决定函数的局部形状。设 \(Du(a)=0\)，而对称矩阵 \(D^2u(a)\) 正定。单位球面紧致，连续函数 \(v\mapsto v^{\mathsf T}D^2u(a)v\) 在其上有正的最小值 \(\lambda\)。当 \(|h|\) 足够小时，二阶余项的模不超过 \(\lambda|h|^2/4\)，于是
\[
 u(a+h)-u(a)\geq\frac{\lambda}{4}|h|^2>0\quad(h\ne0).
\]
因此 \(a\) 是严格局部极小点。负定时对 \(-u\) 使用同一论证；若二次型在两个单位向量上分别为正和负，沿对应的两条直线取充分小的增量，就得到比 \(u(a)\) 大和小的函数值。

\ProseParagraphBreak

二次型退化时，低阶模型可能不能作出判断。三个函数 \(x^4+y^4\)、\(-x^4-y^4\)、\(x^4-y^4\) 在原点的梯度与 Hesse 矩阵都为零，却依次有严格极小点、严格极大点和鞍点。二阶余项虽然比 \(|h|^2\) 小，仍可能决定全部局部行为。使用 Taylor 公式时，先看保留下来的主项是否给出一个严格的符号界，再判断余项是否小到足以维持这个界。

\subsection{逆函数与隐函数定理}
\label{b2:subsec:010103}

若开集 \(U,V\subset\mathbb R^d\) 之间的双射及其逆映射都为 \(C^1\)，则称它为 \(C^1\) \term[weifen-tongpei]{微分同胚}[diffeomorphism]。链式法则说明，这时微分必须可逆。\term[ni-hanshu-dingli]{逆函数定理}[inverse function theorem]在 \(C^1\) 条件下给出局部的逆向结论，也称反函数定理。\termsee[fan-hanshu-dingli]{反函数定理}{逆函数定理}

\begin{theorem}[逆函数]\label{b2:thm:01-inverse}
设 \(O\subseteq\mathbb R^d\) 为开集，\(F:O\rightarrow\mathbb R^d\) 为 \(C^1\) 映射，且 \(a\in O\)。若 \(DF(a)\) 可逆，则存在开邻域 \(O_0\ni a\)、\(V_0\ni F(a)\)，使 \(F:O_0\rightarrow V_0\) 为 \(C^1\) 微分同胚，并且
\[
 D(F^{-1})(y)=\Bigl(DF\bigl(F^{-1}(y)\bigr)\Bigr)^{-1}.
\]
\end{theorem}

第一卷引理 18.1.2 用压缩迭代证明这一结果，其中 \(DF\) 在 \(a\) 的连续性给出一致压缩估计。

\ProseParagraphBreak

局部结论不能自动拼成全局一一对应。例如
\[
 F(s,t)=\bigl(\mathrm{e}^s\cos t,\mathrm{e}^s\sin t\bigr)
\]
处处具有 Jacobi 行列式 \(\mathrm{e}^{2s}>0\)，但 \(F(s,t+2\pi)=F(s,t)\)。用这种映射作坐标时，必须先限制角变量的范围。

\ProseParagraphBreak

方程 \(F(x,y)=0\) 往往已经把 \(y\) 确定为 \(x\) 的函数，只是没有显式解出。\term[yin-hanshu-dingli]{隐函数定理}[implicit function theorem]通过保留 \(x\) 坐标、只把 \(y\) 换成 \(F\) 的值来处理这一问题。

\begin{theorem}[隐函数]\label{b2:thm:01-implicit}
设 \(F\) 是 \(\mathbb R^p\times\mathbb R^q\) 中开集上的 \(C^1\) 映射，值域为 \(\mathbb R^q\)。若 \(F(a,b)=0\)，且对 \(y\) 变量的偏微分矩阵 \(D_yF(a,b)\) 可逆，则存在开邻域 \(X\ni a\)、\(Y\ni b\) 和唯一的 \(C^1\) 映射 \(\varphi:X\rightarrow Y\)，使
\[
 F(x,y)=0\quad\Longleftrightarrow\quad y=\varphi(x)
 \quad\bigl((x,y)\in X\times Y\bigr).
\]
其微分满足
\[
 D\varphi(x)=-\Bigl(D_yF\bigl(x,\varphi(x)\bigr)\Bigr)^{-1}
 D_xF\bigl(x,\varphi(x)\bigr).
\]
\end{theorem}
\begin{proof}
令 \(G(x,y)=(x,F(x,y))\)。其微分为块矩阵
\[
 DG(a,b)=
 \begin{pmatrix}I_p&0\\D_xF(a,b)&D_yF(a,b)\end{pmatrix},
\]
故可逆。由\cref{b2:thm:01-inverse}，\(G\) 在 \((a,b)\) 的某开邻域 \(O_0\) 上具有 \(C^1\) 逆映射。该逆映射的第一组分量必保持不变，因此写作 \(G^{-1}(x,z)=(x,\psi(x,z))\)。

先在 \(O_0\) 内选择包含 \((a,b)\) 的乘积邻域 \(X_1\times Y\)。再缩小 \(a\) 的邻域为 \(X\subseteq X_1\)，使 \((x,0)\) 都属于逆映射的定义域，且 \(\psi(x,0)\in Y\)；这由开性及 \(\psi(a,0)=b\) 处的连续性保证。令 \(\varphi(x)=\psi(x,0)\)，就有 \(F\bigl(x,\varphi(x)\bigr)=0\)。反之，若 \((x,y)\in X\times Y\) 且 \(F(x,y)=0\)，则 \(G(x,y)=(x,0)\)，由 \(G\) 在 \(O_0\) 上的单射性得到 \(y=\varphi(x)\)。

映射 \(\varphi\) 为 \(C^1\)。在 \(O_0\) 内，\(DG\) 由逆映射的链式法则处处可逆，其块形式保证 \(D_yF\) 可逆。对 \(F\bigl(x,\varphi(x)\bigr)=0\) 求导，得到 \(D_xF+D_yF\circ D\varphi=0\)，遂得所述公式。
\end{proof}

例如单位圆的方程 \(x^2+y^2=1\) 在 \((0,1)\) 附近给出 \(y=\sqrt{1-x^2}\)，导数为 \(-x/y\)。在 \((1,0)\) 处，\(y\) 方向的导数退化，但可改解为 \(x\) 关于 \(y\) 的函数。至于 \(F(x,y)=xy\) 在原点的零集合，它含两条相交直线，任意原点邻域内都不能把全部零集合写成一个 \(y=\varphi(x)\) 的图像。可逆性条件既选择可求解的变量，也排除了这种局部交叉。

\subsection{梯度、散度与 Laplace 算子}
\label{b2:subsec:010104}

\begin{definition}\label{b2:def:01-vector-operators}
设 \(U\subseteq\mathbb R^d\) 为开集，\(u\in C^1(U;\mathbb R)\)。其\term[tidu]{梯度}[gradient]为列向量
\[
 \nabla u\coloneqq(\partial_1u,\ldots,\partial_du)^{\mathsf T}.
\]
对 \(C^1\) 向量场 \(X=(X_1,\ldots,X_d)\)，定义其\term[sandu]{散度}[divergence]为
\[
 \operatorname{div}X\coloneqq\sum_{j=1}^d\partial_jX_j.
\]
对 \(u\in C^2(U)\)，定义\term[laplace-suanzi]{Laplace 算子}[Laplace operator]为
\[
 \Delta u\coloneqq\operatorname{div}(\nabla u)
 =\sum_{j=1}^d\partial_j^2u=\operatorname{tr}(D^2u).
\]
\end{definition}
\symidx{\(\nabla u\)：实值函数的梯度}
\symidx{\(\operatorname{div}X\)：向量场的散度}
\symidx{\(\Delta u\)：Laplace 算子，取各坐标二阶偏导数之和}

梯度把微分写成内积：\(D_vu=\nabla u\cdot v\)。由 Cauchy--Schwarz 不等式，单位方向上的变化率不超过 \(\lvert\nabla u\rvert\)，当梯度非零时，在方向 \(\nabla u/\lvert\nabla u\rvert\) 取到最大值。散度则是向量场微分的迹；\secref{b2:sec:0102}的散度定理将赋予它单位体积净流出的积分解释。本书固定 \(\Delta=\sum_j\partial_j^2\) 的符号，后文的正能量通常对应 \(-\Delta\)。

\begin{proposition}\label{b2:prop:01-product-rules}
设 \(U\subseteq\mathbb R^d\) 为开集。对实值 \(u\in C^1(U)\)、\(X\in C^1(U;\mathbb R^d)\)，有
\[
 \operatorname{div}(uX)=\nabla u\cdot X+u\cdot\operatorname{div}X.
\]
若 \(u,v\in C^2(U)\)，则
\[
 \Delta(uv)=u\cdot\Delta v+2\nabla u\cdot\nabla v+v\cdot\Delta u.
\]
\end{proposition}
\begin{proof}
对第 \(j\) 个分量使用一元乘积法则，得到
\(\partial_j(uX_j)=(\partial_ju)X_j+u\cdot\partial_jX_j\)，求和即得第一式。第二式由
\(\partial_j^2(uv)=u\cdot\partial_j^2v+2(\partial_ju)\cdot(\partial_jv)+v\cdot\partial_j^2u\)
对 \(j\) 求和得到。
\end{proof}

虽然 \(\Delta\) 写成坐标偏导数之和，它不依赖欧氏空间中标准正交坐标系的选择。设 \(Q\) 为正交矩阵，在 \(Qz\in U\) 的开集上令 \(v(z)=u(Qz)\)。两次应用链式法则得到
\[
 \nabla v(z)=Q^{\mathsf T}\nabla u(Qz),\quad
 D^2v(z)=Q^{\mathsf T}D^2u(Qz)Q.
\]
取矩阵迹，并利用 \(QQ^{\mathsf T}=I_d\)，即得 \(\Delta v(z)=(\Delta u)(Qz)\)。一般可逆线性变换不满足这个结论：例如把一个坐标放大两倍，会把该方向的二阶导数放大四倍。因而换坐标后，必须重新计算二阶项的系数。

\ProseParagraphBreak

这些算子也适合利用对称性计算。若 \(u(x)=g(r)\)，其中 \(r=\lvert x\rvert>0\)、\(g\in C^2((0,\infty))\)，则链式法则给出
\[
 \partial_ju=g'(r)\frac{x_j}{r},\quad
 \partial_j^2u=g''(r)\frac{x_j^2}{r^2}
 +g'(r)\left(\frac1r-\frac{x_j^2}{r^3}\right).
\]
利用 \(\sum_jx_j^2=r^2\)，得到
\begin{equation}\label{b2:eq:01-radial-laplacian}
 \Delta u=g''(r)+\frac{d-1}{r}\,g'(r).
\end{equation}
例如 \(\Delta\lvert x\rvert^2=2d\)；当 \(d\geq3\) 时，\(\Delta\lvert x\rvert^{2-d}=0\)，而 \(d=2\) 时 \(\Delta\log\lvert x\rvert=0\)。后两式只在原点以外成立，奇点不能由一次代入计算抹去。径向对称性把多元算子化成一元微分表达式，原点处留下的问题则需要后续的积分与弱解方法。

% !TeX root = ../../Book2.tex
\section{多重积分与向量分析}
\label{b2:sec:0102}

偏微分方程的积分计算常有两种目的：选取适合区域的坐标，或者把区域内部的导数转移到边界。前者由换元公式完成，后者由散度与 Stokes 公式完成。本节固定面积元和方向约定，并把这些公式整理为后面分部积分所需的形式。积分均取 Lebesgue 意义；对连续函数在紧致光滑对象上的积分，它与通常的微积分定义一致。

\subsection{变量替换}
\label{b2:subsec:010201}

设 \(U,V\subseteq\mathbb R^d\) 为开集，\(\Phi:U\rightarrow V\) 是 \(C^1\) 微分同胚。其导数的行列式为 Jacobi 行列式。体积伸缩使用它的绝对值，记作 \(J_\Phi\coloneqq|\det D\Phi|\)。第一卷定理 7.5.7 给出如下\term[huanyuan-gongshi]{换元公式}[change of variables formula]：对 Lebesgue 可测集 \(E\subseteq U\) 和 Lebesgue 可测函数 \(f:V\to[0,\infty]\)，
\[
 \int_{\Phi(E)}f(x)\,\mathrm dx
 =\int_E f\bigl(\Phi(y)\bigr)\cdot J_\Phi(y)\,\mathrm dy.
\]
两边允许无穷。对实值或复值函数，先对 \(|f|\) 应用此式检查绝对可积性，再对正负部或实虚部分别使用公式。行列式的符号描述定向，绝对值描述体积；反射坐标不会产生负体积。

\ProseParagraphBreak

例如椭圆内部 \(E=\{\,x^2/a^2+y^2/b^2<1\,\}\)，其中 \(a,b>0\)，适合取
\[
 x=ar\cdot\cos\theta,\quad y=br\cdot\sin\theta,
 \quad 0<r<1,\quad 0<\theta<2\pi.
\]
该映射在所写开矩形上单射，Jacobi 行列式为 \(abr\)。像遗漏的径向线段是二维零测集，故
\[
 \int_E1\,\mathrm dx\,\mathrm dy=\pi ab,
 \quad
 \int_E x^2\,\mathrm dx\,\mathrm dy
 =a^3b\int_0^1r^3\,\mathrm dr\int_0^{2\pi}\cos^2\theta\,\mathrm d\theta
 =\frac{\pi a^3b}{4}.
\]
若把角度扩成 \((0,4\pi)\)，几乎每点被覆盖两次，所得参数积分也加倍。导数处处可逆只保证局部可逆，不能替代这里的全局单射性。

\subsection{曲面积分}
\label{b2:subsec:010202}

边界通常是 \((d-1)\) 维对象，不能用 \(d\) 维体积测量。先看三维空间中最常见的图面：设 \(U\subset\mathbb R^2\) 开，\(g:U\to\mathbb R\) 为 \(C^1\) 函数，以
\[
 X(s,t)=(s,t,g(s,t))
\]
描述它。参数方向上的两个切向量为 \(X_s=(1,0,g_s)\)、\(X_t=(0,1,g_t)\)。一个很小的参数矩形经过线性化后，变成由这两个向量张成的平行四边形，其面积伸缩为
\[
 |X_s\times X_t|=|(-g_s,-g_t,1)|
 =\sqrt{1+g_s^2+g_t^2}.
\]
因此图面面积元的计算公式应当是
\[
 \mathrm dS=\sqrt{1+|\nabla g(s,t)|^2}\,\mathrm ds\,\mathrm dt.
\]
例如对实常数 \(\alpha,\beta\)，斜平面 \(g(s,t)=\alpha s+\beta t\) 上的一块面积，等于它在水平面上的投影面积乘以 \(\sqrt{1+\alpha^2+\beta^2}\)。这个因子测量倾斜带来的伸长，并非取有符号的投影。下面将这一切平面的计算写成适用于一般坐标的形式，再给出面积的坐标无关定义及公式成立的依据。

\ProseParagraphBreak

对 \(1\leq k\leq d\)，若一个集合局部能用定义于开集 \(U\subset\mathbb R^k\) 的单射 \(C^1\) 参数图 \(X:U\rightarrow\mathbb R^d\) 描述，\(X\) 到其相对开像是同胚、\(DX\) 的秩恒为 \(k\)，且参数图之间的过渡映射是 \(C^1\) 微分同胚，则称其为 \term[c1-qumian]{\(C^1\) 曲面}[\(C^1\) surface]。这里维数为 \(k\)，不限定为二维。记其 \term[kwei-jacobi-yinzi]{\(k\) 维 Jacobi 因子}[k-dimensional Jacobian]为
\[
 J_kX\coloneqq\sqrt{\det\bigl((DX)^{\mathsf T}DX\bigr)}.
\]
这是参数小块在切平面中的 \(k\) 维体积伸缩因子。矩阵 \((DX)^{\mathsf T}DX\) 的第 \((i,j)\) 项是两个切向量 \(\partial_iX,\partial_jX\) 的内积，称为 Gram 矩阵；其行列式给出这些向量张成的平行多面体的体积平方。满秩条件保证这个数为正。在三维的二维情形，它就是上面用过的 \(|X_s\times X_t|\)；在平面曲线情形，它是参数速度 \(|X'|\)。环境维数大于参数维数时，\(DX\) 不是方阵，不能直接对它取行列式，也不能把曲面的环境体积为零误认为曲面积为零。

\begin{definition}\label{b2:def:01-vector-surface-integral}
设 \(S\subseteq\mathbb R^d\) 为 \(k\) 维 \(C^1\) 曲面，\(1\leq k\leq d\)。对这个曲面，沿用第一卷定义 15.1.2 的\term[hausdorff-cedu]{Hausdorff 测度}[Hausdorff measure] \(\mathcal H^k\)，其归一化使它在 \(k\) 维仿射平面上等于 \(k\) 维 Lebesgue 测度。以限制 \(\mathcal H^k\llcorner S\) 作为面积测度，在一般维数公式中简写其积分元为 \(\mathrm dS\)。可测函数 \(h:S\rightarrow\mathbb R\) 的\term[qumian-jifen]{曲面积分}[surface integral]定义为
\[
 \int_Sh\,\mathrm dS\coloneqq\int_Sh\,\mathrm d\mathcal H^k,
\]
其中要求 \(h\geq0\)，或 \(\int_S|h|\,\mathrm d\mathcal H^k<\infty\)。复值函数按实部和虚部分别积分，并要求绝对可积。
\symidx{\(\mathrm dS\)：曲面面积元}
\end{definition}

本节凡同时出现曲面和它的边界曲线时，二维曲面的面积元写作 \(\mathrm dS\)，一维曲线的弧长元改写作 \(\mathrm d\ell\)。两者分别是相应对象上的 \(\mathcal H^2\) 与 \(\mathcal H^1\)，从记号即可辨认积分的维数；这不改变一般维数公式中的测度约定。
\symidx{\(\mathrm d\ell\)：曲线弧长元}

第一卷推论 17.3.11（局部面积公式）的单射情形说明，对一张参数图及可测参数集 \(A\subseteq U\)，有
\begin{equation}\label{b2:eq:01-vector-surface-param}
 \int_{X(A)}h\,\mathrm dS
 =\int_A h\bigl(X(s)\bigr)\cdot J_kX(s)\,\mathrm ds.
\end{equation}
这也证明参数表达不依赖坐标选择。事实上，若 \(Y=X\circ\psi\) 是另一张参数图，链式法则与行列式乘法给出 \(J_kY=(J_kX)\circ\psi\cdot|\det D\psi|\)，再在参数域作一次换元即可回到\eqref{b2:eq:01-vector-surface-param}。需要多张图时，先取可数图覆盖 \(S_j\)，再令 \(E_1=S_1\)、\(E_j=S_j\setminus\bigcup_{i<j}S_i\)，在每个不交的 \(E_j\) 上计算并相加。这样的可数覆盖由 \(\mathbb R^d\) 的可数拓扑基取得；各图的相对开像和 \(E_j\) 都是可测的。非负情形用可列可加性，可积情形用绝对收敛，便不会把重叠部分重复计入。曲面上若有边界，边界图可在半空间上作相同计算；光滑边界在 \(k\) 维面积中为零，但在边界自身的 \((k-1)\) 维积分中必须保留。

\ProseParagraphBreak

对开头的图面，\eqref{b2:eq:01-vector-surface-param}现在严格给出了 \(\mathrm dS=\sqrt{1+|\nabla g|^2}\,\mathrm ds\,\mathrm dt\)。在三维空间的二维曲面上，选定连续单位法向 \(n\) 后，向量场 \(F\) 沿该方向的\term[tongliang]{通量}[flux]为 \(\int_SF\cdot n\,\mathrm dS\)，要求被积函数绝对可积。若 \(n=(X_s\times X_t)/|X_s\times X_t|\)，通量就是 \(\int (F\circ X)\cdot(X_s\times X_t)\,\mathrm ds\,\mathrm dt\)。反转 \(n\) 改变通量的符号，不改变标量曲面积分。

\ProseParagraphBreak

取单位球面的上半部 \(S\)，选择向上的法向，并令 \(F=(0,0,1)\)。用 \(g(s,t)=\sqrt{1-s^2-t^2}\) 表示开上半球，便有
\[
 \int_SF\cdot n\,\mathrm dS
 =\int_{s^2+t^2<1}1\,\mathrm ds\,\mathrm dt=\pi.
\]
这一通量是球面在水平面上的投影面积，曲面自身的面积则为 \(2\pi\)。赤道是二维面积零集，加入它不改变两种积分。图函数的导数在靠近赤道时无界，计算仍合法：先在较小同心圆盘上应用局部面积公式，再对非负积分取递增极限。无需假定整张图具有统一的导数界。

\subsection{Gauss 散度定理}
\label{b2:subsec:010203}

设 \(\Omega\subseteq\mathbb R^d\) 为开集。若每个边界点附近，经过平移与正交坐标变换，\(\Omega\) 都是一张 \(C^1\) 函数图像的一侧，则称其具有 \term[c1-bianjie]{\(C^1\) 边界}[\(C^1\) boundary]。这里不要求 \(\Omega\) 连通。\term[wai-danwei-faxiang]{外单位法向}[outward unit normal]指垂直于边界切平面、指向区域外侧的单位向量，记为 \(\nu\)；例如 \(\Omega=\{\,x_d<g(x')\,\}\) 的图边界上，
\[
 \nu=\frac{(-\nabla g,1)}{\sqrt{1+|\nabla g|^2}}.
\]
单侧条件不能略去：从空间删去一张超平面后，区域在该平面的两侧都有点，因而不存在上述意义的外侧。

\begin{theorem}[Gauss 散度定理]\label{b2:thm:01-vector-gauss}
设 \(d\geq1\)，\(\Omega\subseteq\mathbb R^d\) 是有界开集且具有 \(C^1\) 边界，\(F\) 是 \(\overline\Omega\) 某开邻域上的实 \(d\) 维 \(C^1\) 向量场，\(\nu\) 是外单位法向，\(\mathrm dS\) 是边界面积测度（一维为计数测度）。则有\term[gauss-sandu-dingli]{Gauss 散度定理}[Gauss divergence theorem]
\begin{equation}\label{b2:eq:01-vector-gauss}
 \int_\Omega\operatorname{div}F\,\mathrm dx
 =\int_{\partial\Omega}F\cdot\nu\,\mathrm dS.
\end{equation}
\end{theorem}
\begin{proof}
当 \(d=1\) 时，有界且具有光滑边界的开集是有限个不交区间的并。对每个区间应用一维微积分基本定理，左端点取法向 \(-1\)、右端点取法向 \(1\)，并以计数测度解释边界积分，即得公式。以下设 \(d\geq2\)。

第一卷命题 27.5.2（光滑区域的周长）的图表证明已经建立：对 \(C^1\) 图边界及紧支集 \(C^1\) 场 \(G\)，域内散度积分等于上述外法向通量。该证明逐坐标使用一维微积分基本定理，再用有限个图表及单位分解相加。

在本定理中，取 \(\overline\Omega\) 的一个邻域上恒等于 \(1\)、支集紧含 \(F\) 定义域的光滑函数 \(\chi\)。把 \(\chi\cdot F\) 在该定义域之外补零，即得可应用前述公式的场 \(G\)。它在 \(\overline\Omega\) 附近等于 \(F\)，两边便是所求积分。有界性与边界紧致性保证有限图覆盖和有限面积，故连续被积函数都可积。
\end{proof}

例如在半径为 \(R\) 的球上取 \(F(x)=x\)，则 \(\operatorname{div}F=d\)、\(F\cdot\nu=R\)，从\eqref{b2:eq:01-vector-gauss}得到
\[
 d\cdot|B(0,R)|=R\cdot\mathcal H^{d-1}\bigl(\partial B(0,R)\bigr).
\]
有孔区域的内边界同样要计入通量；它的外法向指向孔内。\cref{b2:fig:01-vector-annulus}同时给出稍后平面 Green 定理所用的切向方向。

\ProseParagraphBreak

散度定理要求场在整个闭区域附近具有指定正则性。若场在区域内部某点无定义，即使散度在其余点恒为零，也不能直接断言外边界总通量为零。应先删去该点附近的小球，对带孔区域应用定理，把新产生的内边界项留下，再分析缩球极限。后面研究基本解时，这个内边界项会记录集中于奇点的源。

% 自拟示意图：圆环的区域外法向及与其相配的边界切向，不取自第三方图片。
\begin{figure}[htbp]
\centering
\begin{tikzpicture}[scale=1.05,>=stealth]
 \fill[AnalysisBlue!10,even odd rule] (0,0) circle (2) (0,0) circle (0.9);
 \draw[thick] (0,0) circle (2);
 \draw[thick] (0,0) circle (0.9);
 \node at (-1.25,1.05) {\(\Omega\)};
 \node at (0.3,0.32) {孔};
 \draw[->,thick,AnalysisBlue] (2,0)--(2.8,0) node[right] {\(\nu\)};
 \draw[->,thick] (2,0)--(2,0.8) node[above] {\(\tau\)};
 \draw[->,thick,AnalysisBlue] (-0.9,0)--(-0.15,0);
 \node[AnalysisBlue] at (-0.48,-0.25) {\(\nu\)};
 \draw[->,thick] (-0.9,0)--(-0.9,0.65);
 \node at (-1.14,0.4) {\(\tau\)};
 \draw[->,thick] (110:2) arc[start angle=110,end angle=150,radius=2];
 \draw[->,thick] (300:0.9) arc[start angle=300,end angle=245,radius=0.9];
\end{tikzpicture}
\caption[圆环的外法向与边界方向]{圆环的外法向与边界方向。外法向都指向区域外侧，切向满足 \(\tau=(-\nu_2,\nu_1)\)；因此外圆逆时针，内圆顺时针。}
\label{b2:fig:01-vector-annulus}
\end{figure}

\subsection{Green 恒等式与平面 Green 定理}
\label{b2:subsec:010204}

设 \(\Omega\subset\mathbb R^d\) 是 \(C^1\) 开域，\(\nu\) 为其外单位法向，\(u\) 在 \(\overline\Omega\) 的邻域内为 \(C^1\) 函数。记 \(\partial_\nu u\coloneqq\nabla u\cdot\nu\)，称为沿外法向的\term[faxiang-daoshu]{法向导数}[normal derivative]。以下先对实值函数陈述；复值函数可以分解为实部、虚部使用，能量计算时则需显式加入共轭。

\begin{corollary}[Green 第一与第二公式]\label{b2:cor:01-vector-green}
设 \(\Omega\subseteq\mathbb R^d\) 有界且具有 \(C^1\) 边界，\(\nu\) 为外单位法向；设实值或复值函数 \(u,v\) 在 \(\overline\Omega\) 的某开邻域内为 \(C^2\) 函数。则\term[green-diyi-gongshi]{Green 第一公式}[Green's first identity]为
\begin{equation}\label{b2:eq:01-vector-green-first}
 \int_\Omega\bigl(\nabla u\cdot\nabla v+u\cdot\Delta v\bigr)\,\mathrm dx
 =\int_{\partial\Omega}u\cdot\partial_\nu v\,\mathrm dS,
\end{equation}
而\term[green-dier-gongshi]{Green 第二公式}[Green's second identity]为
\[
 \int_\Omega\bigl(u\cdot\Delta v-v\cdot\Delta u\bigr)\,\mathrm dx
 =\int_{\partial\Omega}
   \bigl(u\cdot\partial_\nu v-v\cdot\partial_\nu u\bigr)\,\mathrm dS.
\]
\end{corollary}
\begin{proof}
对 \(F=u\cdot\nabla v\) 使用散度定理，并代入乘积规则
\(\operatorname{div}(u\cdot\nabla v)=\nabla u\cdot\nabla v+u\cdot\Delta v\)，得到第一式。交换 \(u,v\) 后两式相减，得到第二式。
\end{proof}

对实值函数取 \(v=u\)，\eqref{b2:eq:01-vector-green-first}把梯度的平方积分与 Laplace 算子联系起来：
\[
 \int_\Omega|\nabla u|^2\,\mathrm dx
 =-\int_\Omega u\cdot\Delta u\,\mathrm dx
   +\int_{\partial\Omega}u\cdot\partial_\nu u\,\mathrm dS.
\]
右侧包含方程中的量和边界数据，左侧非负，因而适合估计解及证明唯一性。对复值 \(u\)，应把第一公式中的两函数取为 \(\overline u,u\)；所得 \(\nabla\overline u\cdot\nabla u=|\nabla u|^2\)，不能省略共轭而把 \(\sum_j(\partial_ju)^2\) 当成非负数。

\ProseParagraphBreak

设 \(\Omega\subset\mathbb R^d\) 开，\(u\in C^2(\Omega)\)、\(\varphi\in C_c^\infty(\Omega)\)，两者实值或复值。另一个重要选择是让测试函数的支集紧含 \(\Omega\)。此时边界项消失，得到
\[
 \int_\Omega(-\Delta u)\cdot\varphi\,\mathrm dx
 =\int_\Omega\nabla u\cdot\nabla\varphi\,\mathrm dx.
\]
这里还可以去掉 \(\Omega\) 本身的边界正则性和有界性：把光滑 \(u\) 乘一个在 \(\operatorname{supp}\varphi\) 附近恒为 \(1\) 的内部截断函数，再在包含其支集的大球上应用公式即可。两边的实际积分只在该紧集附近发生，所以不再接触原区域的边界。这种转移导数的方式将在弱解的定义中保留，而对 \(u\) 的光滑要求将相应降低。

\ProseParagraphBreak

Green 第一、第二恒等式把标量函数的导数在区域内部与边界之间转移；平面 Green 定理则把二维向量场的旋度积分写成边界环流。它们都从散度定理推出，但被积对象及使用目的不同。为得到后一种形式，设 \(D\subseteq\mathbb R^2\) 为有界 \(C^1\) 区域，\(P,Q\) 在 \(\overline D\) 附近为 \(C^1\)。在每个边界分支上取单位切向 \(\tau=(-\nu_2,\nu_1)\)，并定义
\[
 \int_{\partial D}P\,\mathrm dx+Q\,\mathrm dy
 \coloneqq\int_{\partial D}(P,Q)\cdot\tau\,\mathrm d\ell.
\]
\begin{theorem}[平面 Green 定理]\label{b2:thm:01-vector-green-plane}
设 \(D\subseteq\mathbb R^2\) 为有界 \(C^1\) 区域，\(P,Q\) 在 \(\overline D\) 的某开邻域上为 \(C^1\) 函数。边界取单位切向 \(\tau=(-\nu_2,\nu_1)\)，其中 \(\nu\) 是外单位法向。则有\term[pingmian-green-gongshi]{平面 Green 定理}[Green's theorem in the plane]
\begin{equation}\label{b2:eq:01-vector-green-plane}
 \int_{\partial D}P\,\mathrm dx+Q\,\mathrm dy
 =\int_D(\partial_xQ-\partial_yP)\,\mathrm dx\,\mathrm dy.
\end{equation}
\end{theorem}
\begin{proof}
对场 \((Q,-P)\) 使用散度定理。其散度为 \(\partial_xQ-\partial_yP\)，法向分量为 \(Q\nu_1-P\nu_2=P\tau_1+Q\tau_2\)，于是得到\eqref{b2:eq:01-vector-green-plane}。
\end{proof}

这个方向使区域位于行进者左侧，故孔的边界须顺时针行进。矩形也满足同式：对 \(D=(a,b)\times(c,d)\)，上下两边的 \(P\) 积分之和为 \(-\int_a^b\int_c^d\partial_yP\,\mathrm dy\,\mathrm dx\)，左右两边的 \(Q\) 积分之和为 \(\int_c^d\int_a^b\partial_xQ\,\mathrm dx\,\mathrm dy\)。一维微积分基本定理与积分次序交换给出所需等式；四个角点不贡献弧长积分。

\subsection{Stokes 公式}
\label{b2:subsec:010205}

在开集 \(U\subset\mathbb R^3\) 中，向量场 \(F=(F_1,F_2,F_3)\in C^1(U;\mathbb R^3)\) 的\term[xuandu]{旋度}[curl]定义为
\[
 \operatorname{curl}F\coloneqq
 (\partial_2F_3-\partial_3F_2,
  \partial_3F_1-\partial_1F_3,
  \partial_1F_2-\partial_2F_1).
\]
\symidx{\(\operatorname{curl}F\)：三维向量场的旋度}
二维曲面上连续选取一个单位法向，就给出了曲面的\term[dingxiang]{定向}[orientation]。带边界 \(C^2\) 曲面在内部由开平面块参数化，在边界附近由闭半平面中的相对开块参数化；边界对应直线部分，参数图及其二阶导数延至该直线附近。沿边界令 \(\eta\) 为曲面切平面中垂直于边界、指向曲面外侧的单位向量，规定边界单位切向
\[
 \tau=n\times\eta.
\]
于是 \(\eta\times\tau=n\)，即“外侧方向在前”的约定。这里 \(\eta\) 仍在曲面内，不能与曲面法向 \(n\) 混淆。沿定向曲线的积分由参数表达 \(\int F(\gamma(t))\cdot\gamma'(t)\,\mathrm dt\) 定义，等价于 \(\int F\cdot\tau\,\mathrm d\ell\)；一维换元保证同向参数变换不改变它。\cref{b2:fig:01-vector-stokes-orientation}将这三个互相垂直的方向放在同一点比较：先确定曲面法向，再找到切平面内的外侧方向，边界切向才随之确定。

% 自拟局部定向示意：只画切平面中的半片和三个方向，不表示曲面的全局形状。
\begin{figure}[htbp]
\centering
\begin{tikzpicture}[x={(1cm,0cm)},y={(0.55cm,0.40cm)},z={(0cm,1cm)},>=stealth]
 \fill[AnalysisBlue!10] (-3,-1,0)--(0,-1,0)--(0,1.5,0)--(-3,1.5,0)--cycle;
 \draw[AnalysisGray] (-3,-1,0)--(0,-1,0);
 \draw[AnalysisGray] (-3,-1,0)--(-3,1.5,0)--(0,1.5,0);
 \draw[thick] (0,-1,0)--(0,1.65,0) node[above right] {\(\partial S\)};
 \node at (-1.7,0,0) {曲面内侧};
 \node at (1.65,-1,0) {外侧};
 \fill (0,0,0) circle (1.5pt);
 \node[below left] at (0,0,0) {\(p\)};
 \draw[->,very thick,AnalysisBlue] (0,0,0)--(0,0,1.65) node[above] {\(n\)};
 \draw[->,very thick] (0,0,0)--(1.8,0,0) node[right] {\(\eta\)};
 \draw[->,very thick] (0,0,0)--(0,1.45,0) node[above left] {\(\tau\)};
\end{tikzpicture}
\caption[曲面定向与诱导边界方向]{曲面定向与诱导边界方向。浅色半片是边界点 \(p\) 处的切平面示意，粗线表示边界的切线方向；\(n\) 垂直于切平面，\(\eta\) 留在切平面内并指向外侧，\(\tau=n\times\eta\) 沿边界。箭头长度不按单位长度绘制。}
\label{b2:fig:01-vector-stokes-orientation}
\end{figure}

每张参数图都能由叉积给出局部法向，但这些选择未必能在整张曲面上相容。定向要求在重叠处一致，沿图表绕行一周后也不改变符号；因此它是一项整体假设。选定定向后可调整参数次序，使每张图都满足 \(X_s\times X_t\) 与 \(n\) 同向。相容图之间的过渡行列式为正，正向的曲线积分才能在不同坐标中接起来。

\ProseParagraphBreak

从一张参数图推广到整张曲面，需要让局部计算在图的人工边缘处自动消失。光滑截断函数将把积分限制在每张图的内部，单位分解再把这些局部积分相加。

\begin{lemma}[紧致曲面的局部化]\label{b2:lem:01-surface-localization}
设 \(S\subset\mathbb R^3\) 为紧致二维 \(C^2\) 曲面，允许有 \(C^2\) 边界，\(W\subset\mathbb R^3\) 为包含 \(S\) 的开集。给定覆盖 \(S\) 的参数图 \(X_\alpha:Q_\alpha\to S\cap O_\alpha\)，其中 \(O_\alpha\subset\mathbb R^3\) 为环境开集；内部图的 \(Q_\alpha\) 是开矩形，边界图的 \(Q_\alpha\) 是半矩形 \((-a_\alpha,a_\alpha)\times[0,b_\alpha)\)，\(a_\alpha,b_\alpha>0\)。边界图的底边对应 \(\partial S\)，参数图延至底边附近为 \(C^2\) 映射。则可选有限个图 \(\alpha_j\)，并找到实函数 \(\rho_j\in C_c^\infty(W\cap O_{\alpha_j})\)，在 \(W\) 中补零后满足
\[
 \rho_j\geq0,\qquad \sum_j\rho_j=1
 \quad\text{于 }S\text{ 的某个环境开邻域}.
\]
每个 \(\rho_j\) 在 \(S\) 上的支集紧含对应参数图；其参数表达在内部矩形的所有边缘附近为零，在边界半矩形的两条侧边及上边附近为零，但不要求在代表 \(\partial S\) 的底边附近为零。
\end{lemma}
\begin{proof}
先在环境空间中构造一个光滑截断。令
\[
 h(t)=\begin{cases}e^{-1/t},&t>0,\\0,&t\leq0,\end{cases}
 \qquad q(t)=\frac{h(t)}{h(t)+h(1-t)}.
\]
\(h\) 在 \(t>0\) 处的各阶导数都是关于 \(1/t\) 的多项式乘 \(e^{-1/t}\)，在 \(t\downarrow0\) 时均趋于零，所以 \(h\) 光滑。分母处处为正，故 \(q\) 光滑，且 \(0\leq q\leq1\)，在 \(t\leq0\) 时为零，在 \(t\geq1\) 时为一。

每个 \(p\in S\) 都能选取 \(r_p>0\)，使闭球 \(\overline{B(p,2r_p)}\) 留在 \(W\) 及指定的环境图邻域中；必要时先缩小图邻域。由紧致性，有限个小球 \(B(p_j,r_j)\) 覆盖 \(S\)。定义
\[
 \chi_j(x)=q\left(\frac{4r_j^2-|x-p_j|^2}{3r_j^2}\right),
 \qquad G(x)=\sum_j\chi_j(x).
\]
于是 \(\chi_j=1\) 于 \(\overline{B(p_j,r_j)}\)，在 \(B(p_j,2r_j)\) 外为零，且 \(G\geq1\) 于 \(S\)。归一化前还须处理 \(G\) 可能为零的位置。令
\[
 R(s)=\begin{cases}q(4s-1)/s,&s>0,\\0,&s\leq0,\end{cases}
 \qquad \rho_j(x)=\chi_j(x)R(G(x)).
\]
由于 \(q(4s-1)=0\) 于 \(s\leq1/4\)，函数 \(R\) 在零点附近恒为零，因而也是光滑函数。每个 \(\rho_j\) 都是非负光滑函数，支集仍留在对应的闭球内，并且
\[
 \sum_j\rho_j=q(4G-1)=1\quad\text{于 }\{G>1/2\}\supset S.
\]
这给出要求的环境单位分解，也使后面使用的 \(\nabla\rho_j\) 有明确含义。

最后，\(\operatorname{supp}\rho_j\cap S\) 是紧含图像中的紧集，其参数原像也紧。内部图的参数域是开矩形，所以该紧集与四条边保持正距离；边界图的底边属于半矩形本身，只有侧边和上边是人工边缘，紧性只要求远离后三条边。因此参数系数可以在人工边缘附近补零，不会给真正的曲面边界加上消失条件。
\end{proof}

\begin{theorem}[Stokes]\label{b2:thm:01-vector-stokes}
设 \(S\subseteq\mathbb R^3\) 是紧致、已定向的二维 \(C^2\) 曲面，允许有 \(C^2\) 边界；\(W\subset\mathbb R^3\) 为包含 \(S\) 的开集，\(F\in C^1(W;\mathbb R^3)\)。取与定向相符的单位法向 \(n\)，沿边界取 \(\tau=n\times\eta\)，其中 \(\eta\) 是曲面内指向外侧、垂直于边界的单位向量，则有\term[stokes-gongshi]{Stokes 公式}[Stokes' theorem]
\begin{equation}\label{b2:eq:01-vector-stokes}
 \int_S\operatorname{curl}F\cdot n\,\mathrm dS
 =\int_{\partial S}F\cdot\tau\,\mathrm d\ell.
\end{equation}
若 \(\partial S=\varnothing\)，右边为零。
\end{theorem}
\begin{proof}
证明分成三个阶段。先建立参数图上的导数恒等式，再核对平面 Green 定理传递的边界方向，最后用\cref{b2:lem:01-surface-localization}把局部等式相加。

第一阶段，在正向参数图 \(X(s,t)\) 上计算。令
\[
 A=(F\circ X)\cdot X_s,\quad B=(F\circ X)\cdot X_t.
\]
乘积规则给出
\[
 \begin{aligned}
 \partial_sB-\partial_tA
 &=\sum_{i,j=1}^3(\partial_iF_j)\circ X
     \cdot\bigl(\partial_sX_i\cdot\partial_tX_j
                 -\partial_tX_i\cdot\partial_sX_j\bigr)\\
 &=(\operatorname{curl}F)\circ X\cdot(X_s\times X_t).
 \end{aligned}
\]
第一行中 \(F\circ X\) 乘混合二阶导数的两项因 \(X_{st}=X_{ts}\) 抵消；把其余六项按 \((i,j)=(1,2),(2,3),(3,1)\) 配对，即为第二行。正向性与面积公式又给出 \(n\,\mathrm dS=(X_s\times X_t)\,\mathrm ds\,\mathrm dt\)。因此曲面旋度通量的被积表达，正是参数平面中的 \(\partial_sB-\partial_tA\)。

第二阶段，取闭包留在参数图内的矩形参数块 \(Q\)，边界图的块允许有一条边落在 \(t=0\) 上；此时使用参数图在底边附近的 \(C^2\) 延拓。因此 \(A,B\) 在闭块附近均为 \(C^1\)，前面直接证明的矩形版 Green 公式给出
\[
 \int_Q(\partial_sB-\partial_tA)\,\mathrm ds\,\mathrm dt
 =\int_{\partial Q}A\,\mathrm ds+B\,\mathrm dt.
\]
对边界参数曲线 \(c\)，链式法则说明右端等于 \(\int F(X(c))\cdot(X\circ c)'\)，即像曲线上的积分。还须核对方向，而不能只核对被积函数。设 \(v=(v_1,v_2)\) 是参数边上的外单位法向，\(w=(-v_2,v_1)\) 是正向单位切向。则
\[
 (DXv)\times(DXw)=\det(v,w)(X_s\times X_t)
 =X_s\times X_t.
\]
\(DXv\) 指向像参数块的外侧，却未必垂直于像边界；将它减去沿 \(DXw\) 的分量，再取单位化，才得到曲面内的外单位向量 \(\eta\)。减去这个切向分量不改变上式叉积，故单位像切向 \(\tau=DXw/|DXw|\) 满足 \(\eta\times\tau=n\)，也就是 \(\tau=n\times\eta\)。参数角点只有有限个，不贡献弧长积分。对于下侧为 \(t=0\)、内部为 \(t>0\) 的半矩形，\(v=(0,-1)\)、\(w=(1,0)\)，所以真正边界沿 \(s\) 增大的方向行进。

第三阶段，在 \(F\) 的定义邻域 \(W\) 中应用\cref{b2:lem:01-surface-localization}，选取有限个 \(\rho_j\)，并在各图上把 \(F\) 换成 \(F_j=\rho_j F\)。内部矩形的参数系数 \(A_j,B_j\) 在全部人工边缘附近为零，可以补零后用第二阶段的公式，边界积分为零。边界半矩形的系数在侧边和上边附近为零，只有代表 \(\partial S\) 的底边可能留下积分。因而人工边界各自不贡献任何项，无需另外选择一套互不重叠的曲面分块。

最后把有限个局部等式相加。由于 \(\sum_j\rho_j=1\) 不仅在 \(S\) 上、而且在它的一个环境邻域内成立，\(\sum_j\nabla\rho_j=0\)，所以
\[
 \sum_j\operatorname{curl}(\rho_j\cdot F)
 =\sum_j\bigl(\nabla\rho_j\times F
                   +\rho_j\cdot\operatorname{curl}F\bigr)
 =\operatorname{curl}F
\]
成立。边界积分也因 \(\sum_j\rho_j=1\) 合为 \(\int_{\partial S}F\cdot\tau\,\mathrm d\ell\)，得到所求公式。若曲面无边界，全部图只留下消失的人工边界项。紧致性和图上连续性保证各个积分有限。
\end{proof}

作为\eqref{b2:eq:01-vector-stokes}的符号核对，取平面单位圆盘、向上法向 \(n=(0,0,1)\) 及 \(F=(-y/2,x/2,0)\)。旋度恒为 \(n\)，面积积分为 \(\pi\)；沿 \(\gamma(t)=(\cos t,\sin t,0)\) 逆时针积分也为 \(\int_0^{2\pi}1/2\,\mathrm dt=\pi\)。反转曲面定向时，边界方向必须同时反转。后面使用分部积分时，正则性决定公式是否可用，方向约定则决定边界项的符号。

% !TeX root = ../../Book2.tex
\section{含参变量积分回顾}
\label{b2:sec:0103}
% 证明工具：第一卷定理5.3.5、5.3.2及7.3.2；下文多参数版本与例子独立推导。

Fourier 分析会反复遇到这样的函数：自变量藏在被积函数中，而积分区域往往没有界。即使每个参数值都给出一个有限积分，也不能立刻断言所得函数连续或可微。这里要把两件事分开：有限区域内的局部变化，以及积分区域远处的尾部。

\subsection{含参变量积分的连续性}
\label{b2:subsec:010301}

设 \((X,\mathcal A,\mu)\) 是测度空间，\(U\subset\mathbb R^m\) 为开集。若对每个 \(t\in U\)，函数 \(x\mapsto f(t,x)\) 可积，则
\[
F(t)\coloneqq\int_X f(t,x)\,\mathrm{d}\mu(x)
\]
是一个\term[hancan-bianliang-jifen]{含参变量积分}[parameter integral]。参数空间和积分空间承担不同的角色：连续、可微等性质针对 \(t\)，可测、可积等性质首先针对 \(x\)。本节允许被积函数取复值。
\termsee[hancanliang-jifen]{含参量积分}{含参变量积分}

\begin{proposition}\label{b2:prop:01-parameter-continuity}
固定 \(t_0\in U\)。设每个 \(f(t,\cdot)\) 可测；存在 \(t_0\) 的邻域 \(V\subset U\)、非负函数 \(g\in L^1(\mu)\) 和一个可测零集 \(N\)，使 \(x\notin N\) 时，\(f(\cdot,x)\) 在 \(t_0\) 连续，且
\[
|f(t,x)|\leq g(x)\quad(t\in V).
\]
则 \(F\) 在 \(V\) 上有定义，且在 \(t_0\) 连续；事实上，\(f(t,\cdot)\) 在 \(L^1(\mu)\) 中趋于 \(f(t_0,\cdot)\)。
\end{proposition}
\begin{proof}
任取 \(t_n\in V\) 且 \(t_n\to t_0\)。共同零集外有 \(f(t_n,x)\to f(t_0,x)\)，并有可积控制 \(g\)。第一卷定理 5.3.5 给出 \(L^1\) 收敛，而
\[
|F(t_n)-F(t_0)|\leq\int_X|f(t_n,x)-f(t_0,x)|\,\mathrm{d}\mu(x).
\]
欧氏空间中的连续性可由序列判定，故结论成立。
\end{proof}
% 跨卷引用：b1:thm:dominated-convergence，Book1/Part01/Chapter05/Section0503.tex。

这里用到的是\term[kongzhi-shoulian-dingli]{控制收敛定理}[dominated convergence theorem]。第一卷已证明其一般测度空间版本；本节把它用于欧氏空间中的参数。条件是局部的：证明 \(t_0\) 附近连续，只需控制这个邻域内的被积函数，不必控制全部参数。

\ProseParagraphBreak

例如，取 \(g\in L^1(\mathbb R^n)\)，令
\[
G(\xi)=\int_{\mathbb R^n}g(x)\cdot e^{-\mathrm{i}x\cdot\xi}\,\mathrm{d}x,
\quad \xi\in\mathbb R^n.
\]
指数因子的模恒为一，所以 \(|G(\xi)|\leq\|g\|_1\)，而 \(|g|\) 给出与 \(\xi\) 无关的控制，故 \(G\) 连续。这正是后面研究 Fourier 变换时最先需要的正则性。若要继续对 \(\xi\) 求导，就会多出 \(x\) 的因子；原先的 \(L^1\) 条件未必还能提供控制。

\subsection{积分号下求导}
\label{b2:subsec:010302}

导数是差商的极限。因此，证明\term[jifenhaoxia-qiudao]{积分号下求导}[differentiation under the integral sign]合法，关键在于找出一个控制差商的可积函数。对多参数而言，还须保证得到的各偏导数组合成连续的全微分。

\begin{theorem}\label{b2:thm:01-parameter-differentiation}
设 \(B\subset\mathbb R^m\) 是开球，\(t_*\in B\)，每个 \(f(t,\cdot)\) 可测且 \(f(t_*,\cdot)\in L^1(\mu)\)。假设存在一个可测零集 \(N\)，使每个 \(x\notin N\) 对应的 \(f(\cdot,x)\) 属于 \(C^1(B)\)，并存在 \(g_j\in L^1(\mu)\)、\(g_j\geq0\)，满足
\[
|\partial_{t_j}f(t,x)|\leq g_j(x)
\quad(t\in B,\ x\notin N,\ 1\leq j\leq m).
\]
则所有 \(f(t,\cdot)\) 可积，\(F\in C^1(B)\)，而且
\begin{equation}\label{b2:eq:01-parameter-derivative}
\partial_{t_j}F(t)=\int_X\partial_{t_j}f(t,x)\,\mathrm{d}\mu(x).
\end{equation}
\end{theorem}
\begin{proof}
可先把全部函数在 \(N\) 上改为零。沿 \(t_*\) 与 \(t\) 间的线段应用一元微积分基本定理，得到
\[
|f(t,x)|\leq |f(t_*,x)|+\sum_{j=1}^m|t_j-t_{*,j}|\cdot g_j(x).
\]
右端可积，所以 \(F(t)\) 有定义。对固定 \(t\) 和 \(j\)，在 \(t+he_j\in B\) 时，同样沿线段积分可得
\[
\frac{f(t+he_j,x)-f(t,x)}{h}
=\int_0^1\partial_{t_j}f(t+she_j,x)\,\mathrm{d}s.
\]
因此差商的模不超过 \(g_j(x)\)。这是对复值函数也成立的积分恒等式，并未使用复值函数的实数中值等式。

取任意非零数列 \(h_k\to0\)，左端是可测函数，逐点极限为 \(\partial_{t_j}f(t,x)\)，所以此导数关于 \(x\) 可测。控制收敛允许把差商极限移入积分，给出\eqref{b2:eq:01-parameter-derivative}。最后，当 \(t\to t_0\) 时，\(\partial_{t_j}f(t,x)\) 连续地趋于 \(\partial_{t_j}f(t_0,x)\)，仍由 \(g_j\) 控制，再用控制收敛可知 \(\partial_{t_j}F\) 连续。\cref{b2:prop:01-c1-criterion}表明 \(F\in C^1(B)\)。
\end{proof}

这一定理可以在每个参数点附近分别使用。同一零集以外的 \(C^1\) 条件，使整条参数线段上的微积分恒等式有意义；“每个参数各自几乎处处可导”本身没有保证这一点。

\ProseParagraphBreak

对上一小节的 \(G\)，若再有 \(|x|g(x)\in L^1(\mathbb R^n)\)，则
\[
\partial_{\xi_j}G(\xi)=\int_{\mathbb R^n}(-\mathrm{i}x_j)g(x)\cdot e^{-\mathrm{i}x\cdot\xi}\,\mathrm{d}x.
\]
控制函数是 \(|x_jg(x)|\)。更高阶求导相应要求更高阶矩可积。这说明积分中的衰减条件怎样转化为参数函数的光滑性；它也提醒我们，形式上容易写出的导数公式仍有各自的适用范围。

\begin{proposition}[Leibniz 积分公式]\label{b2:prop:01-leibniz-integral}
设 \(J\) 是开区间，\(a,b\in C^1(J)\)，且 \(a(t)<b(t)\)。设 \(f\) 定义在包含全部 \((t,x)\)、\(a(t)\leq x\leq b(t)\) 的开集上，\(f\) 与 \(\partial_t f\) 连续。则
\[
\frac{\mathrm{d}}{\mathrm{d}t}\int_{a(t)}^{b(t)}f(t,x)\,\mathrm{d}x
=\int_{a(t)}^{b(t)}\partial_t f(t,x)\,\mathrm{d}x
+f\bigl(t,b(t)\bigr)\cdot b'(t)-f\bigl(t,a(t)\bigr)\cdot a'(t).
\]
\end{proposition}
\begin{proof}
固定 \(t_0\)，取包含 \([a(t_0),b(t_0)]\) 的稍大紧区间 \([c,d]\)。缩小 \(t_0\) 的邻域后，矩形上 \(f\) 与 \(\partial_t f\) 均连续。令
\[
H(t,s)=\int_c^s f(t,x)\,\mathrm{d}x.
\]
固定 \(s\) 时，\cref{b2:thm:01-parameter-differentiation}给出 \(\partial_t H=\int_c^s\partial_t f\,\mathrm{d}x\)；微积分基本定理给出 \(\partial_s H=f(t,s)\)。两者都连续，故对 \(H(t,b(t))-H(t,a(t))\) 使用链式法则即得结论。\(t_0\) 任意，公式在整个 \(J\) 上成立。
\end{proof}

这个结论称为\term[leibniz-jifen-gongshi]{Leibniz 积分公式}[Leibniz integral rule]。两个端点项来自积分区域的变化。例如
\[
\frac{\mathrm{d}}{\mathrm{d}t}\int_0^t e^{tx}\,\mathrm{d}x
=e^{t^2}+\int_0^t xe^{tx}\,\mathrm{d}x\quad(t>0).
\]
若只对被积函数求导，端点 \(x=t\) 带来的那一项就会遗漏。

\subsection{无穷区间积分}
\label{b2:subsec:010303}

在无穷区间上，“积分存在”还需要说明采用哪一种积分。Lebesgue 可积要求绝对值的积分有限；\term[fanchang-jifen]{反常积分}[improper integral]
\[
\int_0^\infty f(t,x)\,\mathrm{d}x
\coloneqq\lim_{R\to\infty}\int_0^R f(t,x)\,\mathrm{d}x
\]
则可能仅靠正负抵消而存在。下面的尾部判据也适用于后一种情形。

\begin{proposition}\label{b2:prop:01-improper-uniform-tail}
设参数集 \(T\subset\mathbb R^m\)，对每个 \(R>0\)，\(F_R(t)=\int_0^R f(t,x)\,\mathrm{d}x\) 在 \(T\) 上连续。若
\[
\lim_{R\to\infty}\ \sup_{t\in T}\ \sup_{S\geq R}
\left|\int_R^S f(t,x)\,\mathrm{d}x\right|=0,
\]
则每个参数对应的反常积分 \(F(t)\) 存在，且 \(F_R\) 在 \(T\) 上一致趋于连续函数 \(F\)。
\end{proposition}
\begin{proof}
给定 \(\varepsilon>0\)，选取 \(R_0\)，使条件中的上确界在 \(R\geq R_0\) 时小于 \(\varepsilon\)。对固定 \(t\)，这正是有限积分 \(F_R(t)\) 在 \(R\to\infty\) 时的 Cauchy 条件，所以极限存在。令 \(S\to\infty\)，得到 \(\sup_{t\in T}|F(t)-F_R(t)|\leq\varepsilon\)，即一致收敛。

对任意 \(t_0\in T\)，先固定这样一个 \(R\)，再由 \(F_R\) 连续使 \(|F_R(t)-F_R(t_0)|<\varepsilon\)。将 \(F(t)-F(t_0)\) 拆为两段尾部误差和中间的有限积分差，便有绝对值不超过 \(3\varepsilon\)。故 \(F\) 连续。
\end{proof}

绝对尾部 \(\sup_t\int_R^\infty|f(t,x)|\,\mathrm{d}x\to0\) 是上述条件的充分条件，但振荡有时能给出更精细的控制。考虑
\[
I(a)=\int_0^\infty e^{-ax}\frac{\sin x}{x}\,\mathrm{d}x,
\quad a\geq0,
\]
其中 \(x=0\) 处用连续延拓解释。当 \(a>0\) 时积分绝对收敛，且在 \(a_0>0\) 附近，参数导数由 \(e^{-a_0x/2}\) 控制。因此
\[
I'(a)=-\int_0^\infty e^{-ax}\sin x\,\mathrm{d}x=-\frac1{1+a^2}.
\]
末一等式可对 \(e^{(-a+\mathrm{i})x}\) 求原函数后取虚部得到。又因 \(|\sin x/x|\leq1\)，有 \(|I(a)|\leq1/a\)，故 \(I(a)\to0\) 当 \(a\to\infty\)。于是
\[
I(a)=\int_a^\infty\frac{\mathrm{d}s}{1+s^2}
=\frac\pi2-\arctan a\quad(a>0).
\]
要让 \(a\downarrow0\)，不能使用上述依赖 \(a_0\) 的控制函数。令 \(h_a(x)=e^{-ax}/x\)。对 \(S\geq R>0\) 分部积分，并利用 \(h_a\) 非负递减，得到
\[
\left|\int_R^S h_a(x)\sin x\,\mathrm{d}x\right|
\leq h_a(R)+h_a(S)+\int_R^S|h_a'(x)|\,\mathrm{d}x
=2h_a(R)\leq\frac2R.
\]
这个估计对全部 \(a\geq0\) 一致成立。有限区间上的积分又连续地依赖于 \(a\)，所以\cref{b2:prop:01-improper-uniform-tail}给出
\[
\int_0^\infty\frac{\sin x}{x}\,\mathrm{d}x=I(0)=\frac\pi2.
\]
此处是反常积分。在每个区间 \([k\pi+\pi/6,k\pi+5\pi/6]\) 上，\(|\sin x|\geq1/2\)；由调和级数发散可知绝对值积分发散，因而不能把这个结果解释成 \(\sin x/x\) 的 Lebesgue 可积性。

\subsection{与第一卷收敛定理的关系}
\label{b2:subsec:010304}

第一卷的积分理论在这里承担三种不同工作。\term[dandiao-shoulian-dingli]{单调收敛定理}[monotone convergence theorem]处理非负递增逼近，允许极限积分为无穷；控制收敛处理有共同可积上界的极限，并给出 \(L^1\) 误差趋零；\term[fubini-dingli]{Fubini 定理}[Fubini's theorem]在绝对可积时交换积分次序。这些结果分别见第一卷定理 5.3.2、定理 5.3.5 和定理 7.3.2。
% 跨卷引用：b1:thm:monotone-convergence、b1:thm:dominated-convergence，Section0503.tex；b1:thm:fubini，Section0703.tex。

\ProseParagraphBreak

例如，若 \(J\) 是有界区间，\((X,\mu)\) 为 \(\sigma\)-有限测度空间，\(q:J\times X\to\mathbb C\) 联合可测，并且 \(|q(s,x)|\leq g(x)\)、\(g\in L^1(\mu)\)，则
\[
\int_J\int_X|q(s,x)|\,\mathrm{d}\mu(x)\,\mathrm{d}s
\leq |J|\cdot\|g\|_1<\infty.
\]
先由第一卷定理 7.3.1，即\term[tonelli-dingli]{Tonelli 定理}[Tonelli\textquotesingle s theorem]，把非负函数 \(|q|\) 的迭代积分识别为乘积积分，再由上述有限性得到乘积空间上的绝对可积。这就验证了交换 \(s\) 与 \(x\) 两重积分所需的绝对可积条件。如果还要把一个极限移入所得积分，则仍须另外核对相应收敛定理的假设。

\ProseParagraphBreak

实际计算时，先明确积分的意义，再写出需要交换的具体操作。对求导，应估计差商；对无穷区域，应控制尾部；对两重积分，应检查联合可测性与绝对可积性。上面的振荡例说明，一致尾部估计可以在没有全局可积控制函数时发挥作用，但它给出的结论仍须按照反常积分来理解。

% 跨卷引用：b1:thm:tonelli，Book1/Part01/Chapter07/Section0703.tex，定理7.3.1。

% !TeX root = ../../Book2.tex
\section{偏微分方程的基本概念}
\label{b2:sec:0104}

前面三节准备了两种表达变化的方式：用导数描述一个点附近的变化，用积分把内部变化与边界上的量联系起来。偏微分方程把这些关系作为未知函数必须满足的条件。一个完整的问题不仅要给出方程，还要说明定义域、附加条件和解的意义。

\subsection{阶数}
\label{b2:subsec:010401}

为统一记号，把 \(\alpha=(\alpha_1,\ldots,\alpha_n)\in\mathbb N_0^n\) 称为\term[duochong-zhibiao]{多重指标}[multi-index]，并记
\[
|\alpha|=\alpha_1+\cdots+\alpha_n,
\quad
\partial^\alpha=\partial_{x_1}^{\alpha_1}\cdots\partial_{x_n}^{\alpha_n}.
\]
这里 \(\mathbb N_0=\{0,1,2,\ldots\}\)。在导数连续到相应阶时，微分次序可以交换。
\symidx{\(\partial^\alpha u\)：多重指标对应的偏导数}

\begin{definition}\label{b2:def:01-pde-order}
设 \(\Omega\subset\mathbb R^n\) 为开集，\(u:\Omega\to\mathbb R\) 为未知函数。含有 \(u\) 的偏导数的关系式
\[
\mathcal F\bigl(x,(\partial^\alpha u(x))_{|\alpha|\leq k}\bigr)=0
\]
称为\term[pianweifen-fangcheng]{偏微分方程}[partial differential equation]，简称 PDE。其中真正出现的最高导数阶数称为方程的\term[jieshu]{阶数}[order]。若最高阶为 \(k\)，则它是 \(k\) 阶方程。
\end{definition}

例如，\(u_t+c u_x=0\) 是一阶方程，\(u_t-\Delta u=f\) 和 \(u_{tt}-\Delta u=0\) 是二阶方程。这里时间 \(t\) 也计作一个自变量，而 \(\Delta\) 只对空间变量作用。热方程的时间导数是一阶，并不妨碍它按最高偏导数阶数被归为二阶方程。

\ProseParagraphBreak

方程必须先化简再判断阶数，不能把彼此抵消的高阶项计入。最高阶系数在某些点为零，则可能造成局部退化。例如 \(xu_{xx}+u_y=0\) 在 \(x\neq0\) 时含有二阶项，而在 \(x=0\) 处这一项消失。这种退化会影响后面的理论，不能仅凭“二阶”二字判断。

\subsection{偏微分方程的线性与非线性类型}
\label{b2:subsec:010402}

分类时应先看最高阶导数怎样出现，再看系数依赖哪些量。以下讨论已经写成明确表达式的标量方程；乘以处处非零的已知函数不改变其解，但不能通过除以可能为零的表达式随意改变方程。

\begin{definition}\label{b2:def:01-pde-linearity}
设 \(\Omega\subset\mathbb R^n\) 开，\(k\ge1\) 为整数，\(u:\Omega\to\mathbb R\) 是未知函数。对含有至多 \(k\) 阶形式导数的实标量偏微分方程，系数与右端均为定义在相应位置或导数变量上的给定实函数，约定如下。
\begin{enumerate}
\item\label{b2:item:01-pde-linear}若可写成
\[
Lu\coloneqq\sum_{|\alpha|\leq k}a_\alpha(x)\partial^\alpha u=f(x),
\]
其中全部系数和右端都是给定的 \(x\) 的函数，则称为\term[xianxing-pianweifen-fangcheng]{线性偏微分方程}[linear partial differential equation]。\(f=0\) 时称为\term[qici-fangcheng]{齐次方程}[homogeneous equation]，否则称为\term[feiqici-fangcheng]{非齐次方程}[inhomogeneous equation]。
\item\label{b2:item:01-pde-semilinear}若最高阶导数线性地出现，且其系数只依赖 \(x\)，即方程写为
\[
\sum_{|\alpha|=k}a_\alpha(x)\partial^\alpha u
+b\bigl(x,(\partial^\beta u)_{|\beta|<k}\bigr)=0,
\]
则称为\term[banxianxing-fangcheng]{半线性方程}[semilinear equation]。
\item\label{b2:item:01-pde-quasilinear}若最高阶导数仍线性地出现，但其系数还可依赖低阶导数，即
\[
\sum_{|\alpha|=k}a_\alpha\bigl(x,(\partial^\beta u)_{|\beta|<k}\bigr)\partial^\alpha u
+b\bigl(x,(\partial^\beta u)_{|\beta|<k}\bigr)=0,
\]
则称为\term[nixianxing-fangcheng]{拟线性方程}[quasilinear equation]。
\item\label{b2:item:01-pde-fullynonlinear}若最高阶导数本身以非线性方式出现，则称为\term[wanquan-feixianxing-fangcheng]{完全非线性方程}[fully nonlinear equation]。
\end{enumerate}
\end{definition}

按这一定义，线性方程属于半线性方程，半线性方程属于拟线性方程。说明一个具体方程的类型时，通常给出其中最精确的一类。对开集 \(\Omega\subset\mathbb R^n\) 上的 \(k\) 阶线性算子 \(L=\sum_{|\alpha|\le k}a_\alpha(x)\partial^\alpha\)，其中 \(k\ge1\) 为整数、\(a_\alpha\) 为给定实值或复值系数，表达式
\[
\sum_{|\alpha|=k}a_\alpha(x)\partial^\alpha
\]
称为其\term[zhubu]{主部}[principal part]；它记录最高阶作用，低阶项仍可能影响解的增长、唯一性及边界条件。

\begin{table}[htbp]
\centering
\caption{按最高阶导数的出现方式区分方程}
\label{b2:tab:01-pde-types}
\begin{tabular}{lll}
\toprule
方程 & 类型 & 判定依据 \\
\midrule
\(u_t-\Delta u=f(x,t)\) & 线性 & 系数和右端均已给定 \\
\(u_t-\Delta u=u^3\) & 半线性 & 非线性只在零阶项 \\
\(u_t+u u_x=0\) & 拟线性 & 一阶导数的系数含 \(u\) \\
\(u_{xx}u_{yy}-u_{xy}^2=f(x,y)\) & 完全非线性 & 二阶导数相乘 \\
\bottomrule
\end{tabular}
\end{table}

\Cref{b2:tab:01-pde-types}中的第三个方程也可以写成 \(u_t+\partial_x(u^2/2)=0\)。在光滑函数上两种写法等价；当导数不再逐点存在时，后一种散度形式会提示一种可保留的积分表达式。解的意义因此与方程的表达方式密切相关。

\subsection{初值问题}
\label{b2:subsec:010403}

把一个自变量解释为时间时，通常在初始时刻给出函数或若干时间导数的值。这样的附加条件称为\term[chuzhi-tiaojian]{初值条件}[initial condition]，方程与初值条件共同构成\term[chuzhi-wenti]{初值问题}[initial value problem]。一阶时间方程通常给定 \(u(\cdot,0)\)，二阶时间方程通常还给定 \(u_t(\cdot,0)\)；这些是常见模型的安排，是否能任意指定仍须由具体方程决定。

\begin{proposition}\label{b2:prop:01-transport-initial}
设 \(c\in\mathbb R\)、\(g\in C^1(\mathbb R)\)。初值问题
\[
u_t+c u_x=0\quad\text{于 }\mathbb R\times(0,\infty),
\quad u(x,0)=g(x)
\]
在 \(C^1(\mathbb R\times[0,\infty))\) 中有唯一解
\[
u(x,t)=g(x-ct).
\]
\end{proposition}
\begin{proof}
链式法则给出 \(u_t=-c g'(x-ct)\)、\(u_x=g'(x-ct)\)，故所给函数满足方程和初值。反过来，若 \(v\) 是一个解，固定 \((x,t)\)，在 \(0\leq s\leq t\) 上考虑 \(v(x-c(t-s),s)\)。其导数为 \(c v_x+v_t=0\)，因而恒定。在两个端点取值可得 \(v(x,t)=v(x-ct,0)=g(x-ct)\)。
\end{proof}

初始图形以速度 \(c\) 平移，函数值沿直线 \(x-ct=\text{常数}\) 保持不变。这些直线给出解怎样从初始数据传到后来的时刻。后面的特征线方法会把这一计算推广到变系数和非线性的一阶方程。

\subsection{边值问题}
\label{b2:subsec:010404}

设 \(\Omega\) 是有界 \(C^1\) 域，\(\nu\) 表示外单位法向。在边界上规定附加条件得到\term[bianzhi-wenti]{边值问题}[boundary value problem]。最常用的三类是
\begin{align*}
u&=g &&\text{于 }\partial\Omega,\\
\partial_\nu u&=h &&\text{于 }\partial\Omega,\\
\partial_\nu u+\sigma u&=h &&\text{于 }\partial\Omega.
\end{align*}
它们依次称为\term[dirichlet-bianjie-tiaojian]{Dirichlet 边界条件}[Dirichlet boundary condition]、\term[neumann-bianjie-tiaojian]{Neumann 边界条件}[Neumann boundary condition]和\term[robin-bianjie-tiaojian]{Robin 边界条件}[Robin boundary condition]。这里 \(g,h,\sigma\) 都是给定的边界函数。前一种指定函数值，第二种指定法向变化率，第三种联系二者。

\ProseParagraphBreak

若实值函数 \(u\) 表示温度，且在边界附近连续可微，常数热导率 \(\kappa>0\) 对应的物理热通量为 \(q=-\kappa\nabla u\)。因此，\(\partial_\nu u=h\) 中的 \(h\) 是外法向导数数据，外向热通量则为
\[
 j\coloneqq q\cdot\nu=-\kappa\partial_\nu u=-\kappa h.
\]
负号表示热量沿温度降低的方向流动：温度沿外法向上升时，热量反而向域内流。若直接给定外向热通量 \(j\)，就应换成 \(h=-j/\kappa\) 再使用本节的 Neumann 写法；本章把常数热导率归一化为 \(\kappa=1\) 时，\(j=-h\)，二者仍相差一个负号。

\begin{proposition}\label{b2:prop:01-poisson-boundary}
设 \(\Omega\subset\mathbb R^n\) 是有界、连通的 \(C^1\) 域，考虑 \(-\Delta u=f\)，并只讨论能在 \(\overline\Omega\) 的某开邻域内成为 \(C^2\) 函数的实值解。
\begin{enumerate}
\item\label{b2:item:01-dirichlet-unique}给定 Dirichlet 数据 \(u=g\) 时，解至多一个。
\item\label{b2:item:01-neumann-compatible}若给定 Neumann 数据 \(\partial_\nu u=h\)，则解存在的必要条件是
\begin{equation}\label{b2:eq:01-neumann-compatibility}
\int_\Omega f\,\mathrm{d}x+\int_{\partial\Omega}h\,\mathrm{d}S=0.
\end{equation}
满足同一数据的两个解至多相差一个常数。
\end{enumerate}
\end{proposition}
\begin{proof}
设 \(u,v\) 满足同一方程和同一边界条件，令 \(w=u-v\)，则 \(\Delta w=0\)。Green 第一公式\eqref{b2:eq:01-vector-green-first}给出
\[
\int_\Omega|\nabla w|^2\,\mathrm{d}x
=\int_{\partial\Omega}w\partial_\nu w\,\mathrm{d}S
-\int_\Omega w\Delta w\,\mathrm{d}x.
\]
在 Dirichlet 情形下 \(w=0\) 于边界，在 Neumann 情形下 \(\partial_\nu w=0\) 于边界，故右端均为零。\(|\nabla w|^2\) 非负连续，积分为零使 \(\nabla w=0\) 处处成立。连通开集内任意两点可用有限段折线相连，沿每段积分便知 \(w\) 为常数。Dirichlet 情形的零边值使该常数为零。

若 \(u\) 满足 Neumann 问题，对方程积分并用\cref{b2:thm:01-vector-gauss}，得到
\[
\int_\Omega f\,\mathrm{d}x=-\int_\Omega\Delta u\,\mathrm{d}x
=-\int_{\partial\Omega}\partial_\nu u\,\mathrm{d}S,
\]
即\eqref{b2:eq:01-neumann-compatibility}。
\end{proof}

相容条件反映了内部源项与边界通量的平衡。在归一化的实值热传导模型中，外向热通量为 \(j=-h\)，所以\eqref{b2:eq:01-neumann-compatibility}等价于 \(\int_\Omega f\,\mathrm dx=\int_{\partial\Omega}j\,\mathrm dS\)：内部产生的总热量等于流出边界的总热量。例如 \(f\equiv1\)、\(h\equiv0\) 时，上式不可能成立，故不存在这样的解。至多一个解和确有一个解是两项不同任务；上述能量计算证明的是唯一性与存在的必要条件，并未给出任意数据的存在定理。

在一维区间 \((0,L)\) 上，可以直接看出相容条件如何参与构造。给定连续的 \(f\)，把左、右端的外法向导数分别记为 \(h_0,h_L\)。因为左端外法向为 \(-1\)，所需条件是 \(u'(0)=-h_0\)、\(u'(L)=h_L\)。对 \(-u''=f\) 积分两次得到
\[
 u'(x)=-h_0-\int_0^x f(s)\,\mathrm ds,\quad
 u(x)=C-h_0x-\int_0^x(x-s)\cdot f(s)\,\mathrm ds.
\]
右端条件恰好要求 \(\int_0^L f+h_0+h_L=0\)。满足它时，任意常数 \(C\) 都给出一个解；不满足时，没有选择 \(C\) 能修正导数条件。例如取 \(f=1\)、\(h_0=0\)、\(h_L=-L\)，解族为 \(C-x^2/2\)。给定两端函数值则能确定两次积分产生的两个常数，得到唯一解。这两种构造解释了为何 Neumann 问题需要总量平衡，而且只能确定到一个常数。

\subsection{经典解、强解与弱解}
\label{b2:subsec:010405}

如果只允许处处可微的函数，一些自然的极限便会落在讨论范围之外。例如连续函数列可以趋于带折角的函数，光滑介质的极限也可能出现分界面。要在这些情形下保留方程，须明确导数和等式各自的意义。

\ProseParagraphBreak

设 \(\Omega\subset\mathbb R^d\) 是开集，\(\alpha\in\mathbb N_0^d\)。先回顾\term[ruo-daoshu]{弱导数}[weak derivative]的积分定义：对 \(u,g\in L^1_{\mathrm{loc}}(\Omega)\)，若
\[
\int_\Omega u\partial^\alpha\varphi\,\mathrm{d}x
=(-1)^{|\alpha|}\int_\Omega g\varphi\,\mathrm{d}x
\quad\bigl(\varphi\in C_c^\infty(\Omega)\bigr),
\]
则称 \(g\) 为 \(u\) 的 \(\alpha\) 阶弱导数。空间 \(C_c^\infty(\Omega)\) 由支集紧含于 \(\Omega\) 的光滑函数组成，这些函数称为\term[ceshi-hanshu]{测试函数}[test function]。\term[sobolev-kongjian]{Sobolev 空间}[Sobolev space] \(W^{k,p}(\Omega)\)，其中 \(k\in\mathbb N_0\)、\(1\le p\le\infty\)、\(\Omega\subset\mathbb R^d\) 开，由函数及所有不超过 \(k\) 阶的弱导数均属于 \(L^p(\Omega)\) 的实值或复值函数组成；下标 \(\mathrm{loc}\) 表示只在每个紧含的开子域上作此要求。这些约定见第一卷定义 20.4.1、定义 21.1.1 和定义 21.1.2，弱导数的唯一性等基本性质也在第一卷相应小节中建立。本节还记 \(H^1=W^{1,2}\)。
% 跨卷引用：b1:def:weak-derivative，Book1/Part04/Chapter20/Section2004.tex；b1:def:sobolev-first-order、b1:def:sobolev-higher-order，Book1/Part04/Chapter21/Section2101.tex。

\symidx{\(C_c^\infty(\Omega)\)：光滑紧支集测试函数空间}
\symidx{\(W^{k,p}(\Omega)\)：Sobolev 空间}
\symidx{\(H^1(\Omega)\)：\(W^{1,2}(\Omega)\)}

\begin{definition}\label{b2:def:01-classical-strong}
设 \(\Omega\subseteq\mathbb R^n\) 为开集，\(k\geq1\) 为整数，讨论实值或复值未知函数的 \(k\) 阶偏微分方程。若 \(u\in C^k(\Omega)\)，方程在每个点成立，则称 \(u\) 为\term[jingdian-jie]{经典解}[classical solution]。若还讨论初值或边值，则另要求相应导数连续延伸到指定部分，并满足那些附加条件。

\ProseParagraphBreak

若所有不超过 \(k\) 阶的弱导数局部属于 \(L^p\)，即 \(u\in W^{k,p}_{\mathrm{loc}}(\Omega)\)，方程中的各项有意义，且等式几乎处处成立，则称它为本节 \(W^{k,p}_{\mathrm{loc}}\) 意义下的\term[qiang-jie]{强解}[strong solution]。这里 \(1\leq p\leq\infty\)，具体问题必须说明所取的 \(p\) 和函数空间。
\end{definition}

“强解”的用法依赖所研究的方程与空间；例如演化方程还可能分别规定时间和空间的正则性。本节给出的约定用于比较经典导数、弱导数及几乎处处成立的方程，不把它当作所有文献统一的定义。

\begin{definition}\label{b2:def:01-divergence-weak}
设 \(n\geq1\) 为整数、\(\Omega\subseteq\mathbb R^n\) 为开集，实矩阵函数 \(A\in L^\infty_{\mathrm{loc}}(\Omega;\mathbb R^{n\times n})\)。函数 \(u,f\) 同取实值或同取复值，且 \(f\in L^2_{\mathrm{loc}}(\Omega)\)。若 \(u\in H^1_{\mathrm{loc}}(\Omega)\) 满足
\begin{equation}\label{b2:eq:01-weak-divergence}
\int_\Omega A\nabla u\cdot\nabla\varphi\,\mathrm{d}x
=\int_\Omega f\varphi\,\mathrm{d}x
\quad\bigl(\varphi\in C_c^\infty(\Omega)\bigr),
\end{equation}
则称 \(u\) 是方程 \(-\operatorname{div}(A\nabla u)=f\) 的\term[ruo-jie]{弱解}[weak solution]。
\end{definition}

所有积分都在测试函数的紧支集附近进行，局部有界的 \(A\) 和局部平方可积的 \(\nabla u,f\) 保证其有限。复值情形仍采用\eqref{b2:eq:01-weak-divergence}中的双线性积分配对；改用共轭测试函数不会改变解类，因为测试函数空间对取共轭封闭。

\ProseParagraphBreak

设进一步有 \(A\in C^1(\Omega;\mathbb R^{n\times n})\)、\(u\in C^2(\Omega)\)、\(f\in C(\Omega)\)，其中 \(u,f\) 同取实值或复值。分部积分说明经典解满足\eqref{b2:eq:01-weak-divergence}。反过来，在这些光滑性条件下若弱等式成立，连续函数 \(r=-\operatorname{div}(A\nabla u)-f\) 与每个测试函数的积分都为零。对实值测试函数分别取实部、虚部，得到
\[
 \int_\Omega\operatorname{Re}r\cdot\varphi\,\mathrm dx=0,
 \quad
 \int_\Omega\operatorname{Im}r\cdot\varphi\,\mathrm dx=0
 \quad\bigl(\varphi\in C_c^\infty(\Omega;\mathbb R)\bigr).
\]
令 \(h\) 分别为 \(\operatorname{Re}r\) 与 \(\operatorname{Im}r\)。倘若 \(h(x_0)>0\)，它在某个紧含于 \(\Omega\) 的小球内有正下界，取该球内非负且不恒为零的实值光滑测试函数，就得到正积分而矛盾；\(h(x_0)<0\) 同理。因此实部、虚部都处处为零，故 \(r=0\) 处处成立。这个证明同时覆盖实值与复值的弱解。

\ProseParagraphBreak

上述定义只表达区域内部的方程，测试函数在边界附近为零，所以它没有包含边界条件。对零 Dirichlet 数据，一个常用做法是进一步要求 \(u\in H_0^1(\Omega)\)，其中 \(H_0^1\) 是 \(C_c^\infty(\Omega)\) 在 \(H^1\) 范数下的闭包，这与第一卷定义 22.2.1 一致。\symidx{\(H_0^1(\Omega)\)：光滑紧支集函数在 \(H^1\) 中的闭包}此时必须把空间条件与内部弱等式同时写出；不能从内部恒等式读出任意边值。

\ProseParagraphBreak

下面三个一维例子区分这些概念。首先，令 \(u(x)=(\max\{x,0\})^2\)。分段积分可验证 \(u'=2\max\{x,0\}\)、\(u''=2\mathbf1_{(0,\infty)}\) 都是弱导数。因此 \(u\in W^{2,\infty}_{\mathrm{loc}}(\mathbb R)\)，是方程 \(u''=2\mathbf1_{(0,\infty)}\) 的强解；但它在零点没有二阶经典导数，所以不是经典解。

\ProseParagraphBreak

其次，\(v(x)=|x|\) 在零点外满足 \(v''=0\)，却不是 \(-v''=0\) 的弱解。事实上 \(v'=\operatorname{sgn}x\) 是它的弱导数，并且
\[
\int_{\mathbb R}v'(x)\varphi'(x)\,\mathrm{d}x=-2\varphi(0).
\]
取 \(\varphi(0)\neq0\) 就否定了弱等式。只把零点删掉再说“方程几乎处处成立”，忽略了导数在该点积累的作用。

\ProseParagraphBreak

最后，在 \((-1,1)\) 上令
\[
a(x)=\begin{cases}1,&x<0,\\2,&x\geq0,\end{cases}
\quad
w(x)=\begin{cases}x,&x<0,\\x/2,&x\geq0.\end{cases}
\]
函数 \(w\) 连续且分段仿射，其弱导数为左右两段的斜率。因此 \(aw'=1\) 几乎处处，故对任意紧支集测试函数，\(\int aw'\varphi'=\int\varphi'=0\)。这给出 \(-(aw')'=0\) 的弱解，尽管 \(w\) 在零点有折角。\Cref{b2:tab:01-interface-flux}把系数、斜率与它们的乘积并列；系数加倍时，斜率恰好减半，通量因而不变。

\begin{table}[htbp]
\centering
\caption[分段系数、斜率与通量]{分段系数、斜率与通量。这里 \(aw'\) 是方程中的数学通量；若 \(a\) 表示热导率，物理热通量为 \(-aw'\)。}
\label{b2:tab:01-interface-flux}
\begin{tabularx}{\textwidth}{@{}p{0.30\textwidth}*{2}{>{\centering\arraybackslash}X}@{}}
\toprule
量 & \(-1<x<0\) & \(0<x<1\) \\
\midrule
系数 \(a\) & \(1\) & \(2\) \\
斜率 \(w'\) & \(1\) & \(\dfrac{1}{2}\) \\
数学通量 \(aw'\) & \(1\) & \(1\) \\
\bottomrule
\end{tabularx}
\end{table}

函数 \(w\) 没有局部可积的二阶弱导数：分段积分给出
\[
-\int_{-1}^1 w'\varphi'\,\mathrm{d}x=-\frac12\varphi(0).
\]
若右端可以写为某个局部可积函数与 \(\varphi\) 的积分，取 \(\varphi_\varepsilon(x)=\eta(x/\varepsilon)\)，其中 \(\eta\in C_c^\infty((-1,1))\)、\(0\leq\eta\leq1\)、\(\eta(0)=1\)，则积分绝对连续性使左述表示趋于零，而 \(-\varphi_\varepsilon(0)/2=-1/2\)，矛盾。弱方程保留的是通量 \(aw'\) 的平衡，并不要求 \(w\) 的所有二阶弱导数都是函数。

\ProseParagraphBreak

上述折角例子可以推广为界面条件的计算。设 \(n\geq2\) 为整数、\(V\subseteq\mathbb R^{n-1}\) 为非空开集，\(\Omega=(-1,1)\times V\)；取正常数 \(a_-,a_+\)，系数矩阵 \(A\) 在 \(x_1<0\) 和 \(x_1>0\) 分别为 \(a_-I_n\)、\(a_+I_n\)，界面本身的取值不影响积分。给定实数 \(b,m_-,m_+\)，取连续的分段仿射函数
\[
 u(x)=\begin{cases}b+m_-x_1,&x_1<0,\\b+m_+x_1,&x_1\geq0.\end{cases}
\]
连续性使一阶弱导数就是两侧的通常梯度，没有界面上的附加项。令 \(\varphi\in C_c^\infty(\Omega)\)，在 \(x_1\) 方向分别积分，再用 Fubini 定理，有
\[
 \int_\Omega A\nabla u\cdot\nabla\varphi\,\mathrm dx
 =(a_-m_--a_+m_+)\int_V\varphi(0,x')\,\mathrm dx'.
\]
所有被积函数都限于一个共同紧集，因而换序合法。这个积分对全部测试函数为零，当且仅当 \(a_-m_-=a_+m_+\)：必要性可取分离变量测试函数，其界面积分不为零；充分性由所写等式直接得到。因此无源弱方程要求通量连续，而非斜率连续。这个计算也是把分部积分用于分段光滑对象时检查内部边界项的方式。

\ProseParagraphBreak

若两侧通量不等，上式右端描述集中在界面 \(\Sigma=\{0\}\times V\) 上的源。为准确表达它，用 \(\delta_\Sigma\) 表示作用在测试函数上的界面积分规则
\[
 \langle\delta_\Sigma,\varphi\rangle\coloneqq
 \int_V\varphi(0,x')\,\mathrm dx'
 \quad\bigl(\varphi\in C_c^\infty(\Omega)\bigr).
\]
于是
\[
 -\operatorname{div}(A\nabla u)
 =(a_-m_--a_+m_+)\delta_\Sigma
\]
在分布意义下成立；这里的“分布”允许右端是一个对测试函数的积分规则，不要求它由普通函数表示。非零的 \(\delta_\Sigma\) 不能由 \(L^2_{\mathrm{loc}}(\Omega)\) 函数表示：取 \(\eta\in C_c^\infty((-1,1))\)、\(\eta(0)=1\)，再取非负且不恒为零的 \(\theta\in C_c^\infty(V)\)，令 \(\varphi_\varepsilon(x)=\eta(x_1/\varepsilon)\theta(x')\)。界面积分恒为 \(\int_V\theta>0\)；任何局部平方可积函数与 \(\varphi_\varepsilon\) 的积分，却因局部可积性及积分绝对连续性而趋于零，二者不可能相同。因此通量不匹配时，这个等式超出了\cref{b2:def:01-divergence-weak}的 \(L^2_{\mathrm{loc}}\) 右端框架；该定义在这个模型中只包含通量匹配的无源情形。删去界面后得到的几乎处处等式不能消除这个源。后面的分布理论将系统处理这种右端。

\subsection{叠加原理}
\label{b2:subsec:010406}

线性方程的一个基本优点，是可以把较复杂的数据拆开处理。\term[diejia-yuanli]{叠加原理}[superposition principle]必须同时考虑方程与附加条件。

\begin{proposition}\label{b2:prop:01-superposition}
设 \(L\) 是线性微分算子，\(B\) 是线性初值或边界算子，且讨论的函数空间对线性组合封闭。若 \(Lu_j=f_j\)、\(Bu_j=g_j\)，则对任意常数 \(c_1,\ldots,c_N\)，
\[
L\Bigl(\sum_{j=1}^N c_ju_j\Bigr)=\sum_{j=1}^N c_jf_j,
\quad
B\Bigl(\sum_{j=1}^N c_ju_j\Bigr)=\sum_{j=1}^N c_jg_j.
\]
特别地，若 \(u_p\) 满足 \(Lu_p=f\)、\(Bu_p=g\)，则同一问题的全部解恰为 \(u_p+v\)，其中 \(Lv=0\)、\(Bv=0\)。
\end{proposition}
\begin{proof}
两个等式直接来自线性。若 \(u\) 与 \(u_p\) 满足同一数据，则线性给出 \(L(u-u_p)=0\)、\(B(u-u_p)=0\)；反过来，把这样的 \(v\) 加到 \(u_p\) 上就恢复原数据。弱解情形下，对测试函数的积分等式也保持线性，故同一论证适用。
\end{proof}

同一非齐次方程的两个解直接相加，通常得到右端加倍的问题。对非线性方程，连这种结论也不成立：\(u(t)=1/(1-t)\) 在 \(t<1\) 满足 \(u_t=u^2\)，但 \(2u\) 的时间导数是 \(2u^2\)，而其平方是 \(4u^2\)。

\ProseParagraphBreak

后面的 Fourier 方法会把一个函数分解为许多振动模式，再对各模式分别求解。有限叠加由线性直接保证；无穷叠加还涉及级数收敛、逐项微分以及初值或边值的恢复，这些正是接下来几章必须建立的内容。

% 跨卷引用：b1:def:sobolev-zero-boundary，Book1/Part04/Chapter22/Section2202.tex，定义22.2.1。

% 本章小结（不计入编号）
\begin{chaptersummary}
全微分把各方向的变化组织为一个统一的线性近似，链式法则让这种近似在坐标变换下保持一致。换元公式还需要全局单射性，边界积分还需要正则性和相容定向；仅检查导数矩阵可逆，或只看一条边界，都可能遗漏关键条件。Green 公式将方程中的二阶导数转化为梯度能量与边界项，是本章通往偏微分方程的主要工具。

\ProseParagraphBreak

含参变量积分中的极限与求导需要相应的估计：共同可积控制允许取极限，导数控制先约束差商，一致尾部控制处理无穷区域。振荡积分的例子表明，条件收敛也可以连续地依赖参数，但必须保留反常积分的意义，不能把正负抵消当作绝对可积。

\ProseParagraphBreak

一个方程问题由方程、定义域、附加条件和解的意义共同确定。经典解逐点满足方程；本节的强解使用可积弱导数和几乎处处等式；散度型弱解通过测试函数表达通量平衡。线性保证有限叠加，至于无穷叠加何时恢复原函数、能否逐项微分，还需要进一步的收敛理论。

\ProseParagraphBreak

回到章首的稳态热传导模型：Green 公式把温度梯度的能量与内部热源、边界通量联系起来；Neumann 相容条件说明总发热量必须等于流出边界的总热量。由于本章将 Neumann 数据记为 \(h=\partial_\nu u\)，实际外向热通量是 \(-h\)，所以平衡写作 \(\int_\Omega f=-\int_{\partial\Omega}h\)。介质分界面上的折角例子进一步说明，弱形式要求两侧通量相匹配，并不要求温度斜率连续。这些概念确定了需要求解的问题；下一章开始的 Fourier 分析则提供构造解、控制无穷叠加的具体方法。
\end{chaptersummary}


\BookExercises

\Cref{b2:ex:01-implicit-global}与\cref{b2:ex:01-product-ratio-change}训练具体计算及其合法性；\cref{b2:ex:01-directional-continuous}与\cref{b2:ex:01-concentration-at-endpoint}检验两个常见的条件边界。\Cref{b2:ex:01-annulus-orientation}和\cref{b2:ex:01-damped-cosine}分别组合几何定向与参数积分方法；\cref{b2:ex:01-pde-classification-weak-wave}把方程分类连接到有折角的输运解。带星号的\cref{b2:ex:01-robin-energy}进一步考察低阶项和边界项怎样影响唯一性；\cref{b2:ex:01-paraboloid-stokes}综合训练曲面面积元、法向和边界环流。

% 第二卷第1章习题初稿；2026-09-15。全部自拟，按本章已建工具独立推导。
% 正文接入时保留章文件中唯一的 \BookExercises，不重复插入习题标题。

% 自拟：把正文“方向导数全为零但不连续”加强为连续反例。
\begin{exercise}\label{b2:ex:01-directional-continuous}
在 \(\mathbb R^2\) 上定义
\[
 f(x,y)=
 \begin{cases}
 \dfrac{x^3y}{x^4+y^2},&(x,y)\ne(0,0),\\
 0,&(x,y)=(0,0).
 \end{cases}
\]
证明 \(f\) 在原点连续，且沿每个方向的方向导数都为零，但在原点并非 Fréchet 可微。这说明在“所有方向导数存在”之外再加连续性，仍不能替代全微分的余项条件。
\end{exercise}
\begin{solution}
由 \(2|x^2y|\leq x^4+y^2\)，有 \(|f(x,y)|\leq |x|/2\)，故 \(f\) 在原点连续。固定方向 \(v=(a,b)\)。若 \(b\ne0\)，则
\[
 \dfrac{f(ta,tb)}{t}
 =\dfrac{t a^3b}{t^2a^4+b^2}\longrightarrow0.
\]
若 \(b=0\)，差商恒为零。因此所有方向导数都为零。

假如全微分存在，它在每个向量 \(v\) 上的值都必须等于方向导数，因而只能是零映射。然而沿曲线 \((x,y)=(s,s^2)\) 有
\[
 f(s,s^2)=\dfrac{s}{2},\quad
 \dfrac{|f(s,s^2)|}{\sqrt{s^2+s^4}}
 =\dfrac{1}{2\sqrt{1+s^2}}\longrightarrow\dfrac12.
\]
余项不是 \(o\bigl(|(x,y)|\bigr)\)，所以全微分不存在。
\end{solution}

% 自拟：局部隐函数拼成全局函数，再提取原点二阶信息。
\begin{exercise}\label{b2:ex:01-implicit-global}
考虑方程 \(z+z^3=xy\)。证明它对每个 \((x,y)\in\mathbb R^2\) 恰有一个实解 \(z=g(x,y)\)，且 \(g\) 是 \(C^2\) 函数。求 \(Dg(0,0)\) 和 \(D^2g(0,0)\)，并证明
\[
 |g(x,y)-xy|\leq |xy|^3.
\]
说明最后一个估计与二阶 Taylor 公式如何相容。
\end{exercise}
\begin{solution}
实函数 \(h(z)=z+z^3\) 严格递增，且在两端分别趋于正、负无穷，故每个实数 \(xy\) 恰有一个原像。对 \(F(x,y,z)=z+z^3-xy\)，有 \(\partial_zF=1+3z^2>0\)。隐函数定理在每一点给出局部 \(C^1\) 解；唯一性保证重叠邻域上的局部解一致，因而得到全局 \(C^1\) 函数。对恒等式求导可得
\[
 g_x=\dfrac{y}{1+3g^2},\quad
 g_y=\dfrac{x}{1+3g^2}.
\]
右侧都是 \(C^1\) 函数，所以 \(g\in C^2\)。原点处 \(g=0\)，上述公式给出 \(Dg(0,0)=0\)。再求一次导数得
\[
 g_{xx}(0,0)=g_{yy}(0,0)=0,\quad g_{xy}(0,0)=1,
 \quad D^2g(0,0)=\begin{pmatrix}0&1\\1&0\end{pmatrix}.
\]

由 \(xy=g(1+g^2)\) 可知 \(|g|\leq|xy|\)，故
\[
 |g-xy|=|g|^3\leq|xy|^3
 \leq\dfrac{(x^2+y^2)^3}{8}.
\]
二阶 Taylor 多项式为 \(xy\)，而所得余项甚至是 \(O\bigl(|(x,y)|^6\bigr)\)，当然也是 \(o\bigl(|(x,y)|^2\bigr)\)。
\end{solution}

% 自拟：用乘积和比值同时拉直四条边界；不同于正文椭圆换元。
\begin{exercise}\label{b2:ex:01-product-ratio-change}
设
\[
 \Omega=\{\,(x,y)\in(0,\infty)^2\mid1<xy<4,\ 1<y/x<9\,\}.
\]
用变换 \(u=xy\)、\(v=y/x\) 分别计算 \(\Omega\) 的面积和积分
\[
 \int_\Omega xy\,\mathrm dx\,\mathrm dy.
\]
求解时核对变换的单射性与 Jacobi 行列式，不只作形式代入。
\end{exercise}
\begin{solution}
在正象限内，逆变换唯一地写成
\[
 x=\sqrt{u/v},\quad y=\sqrt{uv},
 \quad (u,v)\in(1,4)\times(1,9).
\]
正、逆变换均为 \(C^1\)，所以这是两个开集间的微分同胚。直接计算正向导数得
\[
 \det\dfrac{\partial(u,v)}{\partial(x,y)}
 =\det\begin{pmatrix}y&x\\-y/x^2&1/x\end{pmatrix}
 =\dfrac{2y}{x}=2v.
\]
故逆变换的 Jacobi 行列式绝对值为 \(1/(2v)\)。换元公式和非负函数的迭代积分给出
\[
 |\Omega|=\int_1^4\int_1^9\dfrac{1}{2v}\,\mathrm dv\,\mathrm du
 =3\log3,
\]
以及
\[
 \int_\Omega xy\,\mathrm dx\,\mathrm dy
 =\int_1^4\int_1^9\dfrac{u}{2v}\,\mathrm dv\,\mathrm du
 =\dfrac{15}{2}\log3.
\]
\end{solution}

% 自拟：同一环域比较通量和环流，训练内边界的方向及奇点排除。
\begin{exercise}\label{b2:ex:01-annulus-orientation}
在穿孔平面上令
\[
 F(x,y)=\dfrac{(x,y)}{x^2+y^2},\quad
 G(x,y)=\dfrac{(-y,x)}{x^2+y^2}.
\]
设 \(0<a<b\)，并令 \(A=\{\,(x,y)\mid a<\sqrt{x^2+y^2}<b\,\}\)。
\begin{enumerate}
 \item\label{b2:item:01-annulus-flux} 计算 \(\operatorname{div}F\) 及 \(F\) 在 \(A\) 的两条边界上的向外通量。
 \item\label{b2:item:01-annulus-circulation} 把 \(A\) 视为取向上法向的平面曲面，计算 \(\partial_xG_2-\partial_yG_1\)，以及 \(G\) 沿两条诱导定向边界的环流。
 \item\label{b2:item:01-annulus-singularity} 解释为何不能把内边界删去，再对整个圆盘套用相应公式并断言外边界积分为零。
\end{enumerate}
\end{exercise}
\begin{solution}
在原点外，
\[
 \partial_xF_1=\dfrac{y^2-x^2}{(x^2+y^2)^2},\quad
 \partial_yF_2=\dfrac{x^2-y^2}{(x^2+y^2)^2},
\]
故散度为零。在半径 \(b\) 的外圆上，\(F\cdot\nu=1/b\)，弧长元为 \(b\,\mathrm d\theta\)，通量为 \(2\pi\)。在半径 \(a\) 的内圆上，环域的外法向指向原点，故 \(F\cdot\nu=-1/a\)，通量为 \(-2\pi\)。总通量为零，与散度积分一致。

由于 \(G=(-F_2,F_1)\)，有 \(\partial_xG_2-\partial_yG_1=\operatorname{div}F=0\)。向上法向在外圆诱导逆时针方向，参数化为 \(\gamma(\theta)=(b\cos\theta,b\sin\theta)\) 时，
\[
 G\bigl(\gamma(\theta)\bigr)\cdot\gamma'(\theta)=1.
\]
故外圆环流为 \(2\pi\)。内圆的诱导方向是顺时针，环流为 \(-2\pi\)，两者相消。

这两个向量场都在原点奇异，不能延伸为整个圆盘邻域上的 \(C^1\) 向量场。公式在每个环域上都适用，但令 \(a\downarrow0\) 时，内边界积分恒为 \(-2\pi\)，不会因圆周长度趋零而消失。
\end{solution}

% 自拟提高题：将正文零 Dirichlet/Neumann 模型改为带反应项与边界项的能量判别。
\begin{optionalexercise}\label{b2:ex:01-robin-energy}
设 \(\Omega\subseteq\mathbb R^d\) 是有界连通开集，边界为 \(C^1\)。实函数 \(c\in C(\overline\Omega)\)、\(\beta\in C(\partial\Omega)\) 均非负。考虑实值经典解满足的问题
\[
 -\Delta u+c\cdot u=f\quad\text{于 }\Omega,\quad
 \partial_\nu u+\beta\cdot u=g\quad\text{于 }\partial\Omega,
\]
其中解在闭包的某个邻域内属于 \(C^2\)。
\begin{enumerate}
 \item\label{b2:item:01-robin-positive} 若 \(\int_\Omega c\,\mathrm dx+\int_{\partial\Omega}\beta\,\mathrm dS>0\)，证明解至多一个。
 \item\label{b2:item:01-robin-zero} 若 \(c=\beta=0\)，证明任意两解之差为常数。
 \item\label{b2:item:01-robin-negative} 现在去掉 \(c\geq0\) 的假设。在区间 \((0,\pi)\) 上取 \(c=-1\)、\(\beta=0\)、\(f=g=0\)，构造彼此相差并非常数的多个解，说明反应项的符号会影响唯一性。
\end{enumerate}
\end{optionalexercise}
\begin{solution}
设 \(w\) 是两解之差。它满足 \(-\Delta w+cw=0\) 和 \(\partial_\nu w=-\beta w\)。把方程乘以 \(w\)，再使用 Green 第一公式，得到
\[
 0=\int_\Omega|\nabla w|^2\,\mathrm dx
 +\int_\Omega c\cdot w^2\,\mathrm dx
 +\int_{\partial\Omega}\beta\cdot w^2\,\mathrm dS.
\]
三项都非负，因而 \(\nabla w=0\) 几乎处处。梯度连续，若它在某点非零，则在一个球上有正下界，与积分为零矛盾；所以 \(\nabla w\) 处处为零。函数在每个充分小的球上为常数；由 \(\Omega\) 的连通性，这些局部常数相等，故 \(w=C\)。

在第一种情形中，把常数代入上式得到
\[
 C^2\left(\int_\Omega c\,\mathrm dx+
 \int_{\partial\Omega}\beta\,\mathrm dS\right)=0,
\]
故 \(C=0\)。若 \(c=\beta=0\)，上面的论证只要求 \(w\) 为常数，而给任何已有解加常数仍是解。

最后，对任意实数 \(A\)，函数 \(u_A(x)=A\cos x\) 满足 \(-u_A''-u_A=0\)，且 \(u_A'(0)=u_A'(\pi)=0\)，所以两端外法向导数均为零。不同 \(A\) 给出的解之差通常不是常数。此时能量恒等式含有负项 \(-\int_0^\pi u_A^2\,\mathrm dx\)，不能再逐项推出它们为零。
\end{solution}

% 自拟：从一个二参数收敛积分中提取对数，而非复述积分号下求导定理。
\begin{exercise}\label{b2:ex:01-damped-cosine}
对 \(a>0\)、\(b\in\mathbb R\)，令
\[
 J(a,b)=\int_0^\infty e^{-ax}\cdot\dfrac{\cos(bx)-1}{x}\,\mathrm dx.
\]
证明积分绝对收敛，验证可以对 \(b\) 在积分号下求导，并计算 \(J(a,b)\)。须指出原点附近和无穷远处分别使用的估计。
\end{exercise}
\begin{solution}
当 \(0<x\leq1\) 时，\(|\cos(bx)-1|\leq b^2x^2/2\)，被积函数的模不超过 \(b^2x/2\)；当 \(x\geq1\) 时，它不超过 \(2e^{-ax}\)。因此积分绝对收敛。在任意参数矩形 \(a\geq a_0>0\)、\(|b|\leq B\) 内，以上两个控制函数也可分别取 \(B^2x/2\) 与 \(2e^{-a_0x}\)。

对 \(b\) 的导数为 \(-e^{-ax}\sin(bx)\)，其绝对值由可积函数 \(e^{-a_0x}\) 控制。积分号下求导定理给出
\[
 \partial_bJ(a,b)=-\int_0^\infty e^{-ax}\sin(bx)\,\mathrm dx.
\]
记右侧不带负号的积分为 \(S\)，并记 \(C=\int_0^\infty e^{-ax}\cos(bx)\,\mathrm dx\)。在有限区间上分部积分，再令上限趋于无穷，指数衰减使边界项消失，得到
\[
 S=\dfrac ba C,\quad C=\dfrac1a-\dfrac ba S.
\]
解得 \(S=b/(a^2+b^2)\)。由于 \(J(a,0)=0\)，对 \(b\) 积分便有
\[
 J(a,b)=-\int_0^b\dfrac{s}{a^2+s^2}\,\mathrm ds
 =-\dfrac12\log\left(1+\dfrac{b^2}{a^2}\right).
\]
\end{solution}

% 自拟：光滑集中族；显式求出不可积的共同包络，排除逐点极限直接换序。
\begin{exercise}\label{b2:ex:01-concentration-at-endpoint}
对 \(x\in(0,1)\) 及 \(t\geq0\)，定义
\[
 f(t,x)=
 \begin{cases}
 \dfrac{t}{x^2}e^{-t/x},&t>0,\\
 0,&t=0.
 \end{cases}
\]
证明：每个固定 \(x\) 处都有 \(f(t,x)\to0\) 当 \(t\downarrow0\)，但相应积分不趋于零。进一步证明，对任意 \(\delta\in(0,1)\)，不存在一个 \(g\in L^1(0,1)\)，能在同一个零集之外对所有 \(0<t<\delta\) 满足 \(f(t,x)\leq g(x)\)。
\end{exercise}
\begin{solution}
固定 \(x>0\) 时，\(t e^{-t/x}/x^2\to0\)。另一方面，以 \(y=t/x\) 换元可得
\[
 \int_0^1\dfrac{t}{x^2}e^{-t/x}\,\mathrm dx
 =\int_t^\infty e^{-y}\,\mathrm dy=e^{-t}\longrightarrow1.
\]
积分在参数端点不连续。假如存在题述控制函数，那么对零集外每个 \(0<x<\delta\)，可取 \(t=x\)，得到
\[
 g(x)\geq f(x,x)=\dfrac{1}{ex}.
\]
右侧在 \((0,\delta)\) 上不可积，与 \(g\in L^1\) 矛盾。这里失败的是对参数邻域共同成立的可积控制；每个单独截面可积并不能提供这种控制。
\end{solution}

% 自拟：类别识别结合连续但有折角的输运解；弱恒等式直接由换元和分部积分验证。
\begin{exercise}\label{b2:ex:01-pde-classification-weak-wave}
\begin{enumerate}
 \item\label{b2:item:01-pde-classification} 判断下列方程的阶数，并按线性、半线性、拟线性、完全非线性中最具体的类别分类：
 \[
 \begin{aligned}
 &u_t+\partial_x(xu)=f(x,t),
 &&u_t-u_{xx}+(u_x)^2=0,\\
 &(1+(u_x)^2)u_{xx}+u_{yy}=0,
 &&e^{u_{xx}}+u_{yy}=0.
 \end{aligned}
 \]
 \item\label{b2:item:01-pde-corner-wave} 下面比较经典导数不存在与方程在弱导数、分布意义下成立的区别。固定实数 \(c\ne0\)，令 \(u(x,t)=|x-ct|\) 于 \(Q=\mathbb R\times(0,\infty)\)。求它的一阶弱导数，证明 \(u\in W^{1,\infty}_{\mathrm{loc}}(Q)\)，并说明它是否是\cref{b2:def:01-classical-strong}意义下的强解。证明它不是方程 \(u_t+c u_x=0\) 的经典 \(C^1\) 解，但满足
 \[
 \int_0^\infty\int_{\mathbb R}
 u\cdot(\varphi_t+c\varphi_x)\,\mathrm dx\,\mathrm dt=0
 \quad\bigl(\varphi\in C_c^\infty(\mathbb R\times(0,\infty))\bigr).
 \]
 证明其初值为 \(|x|\)，这里初值按每个紧区间上的一致收敛理解。
\end{enumerate}
\end{exercise}
\begin{solution}
四个方程依次为一阶线性、二阶半线性、二阶拟线性、二阶完全非线性方程。第一个方程展开为 \(u_t+xu_x+u=f(x,t)\)，全部系数已知。第二个方程的最高阶导数 \(u_{xx}\) 线性出现，系数固定；非线性发生在低一阶的导数 \(u_x\) 上。第三个方程的二阶导数仍线性出现，但 \(u_{xx}\) 的系数依赖一阶导数 \(u_x\)，所以它不是半线性的。第四个方程把二阶导数代入指数函数，不能写成最高阶导数的线性组合。

在直线 \(x=ct\) 上，函数关于 \(x\) 的左右导数分别为 \(-1\) 和 \(1\)，所以它不属于 \(C^1\)。在这条直线之外，直接计算仍有 \(u_t+c u_x=0\)。为了检验弱恒等式，不把这一几乎处处的计算直接当成结论，而作换元 \(r=x-ct\)：
\[
 \begin{split}
 &\int_0^\infty\int_{\mathbb R}|x-ct|\cdot
 (\varphi_t+c\varphi_x)(x,t)\,\mathrm dx\,\mathrm dt\\
 &\quad=\int_{\mathbb R}|r|\left(
 \int_0^\infty\dfrac{\mathrm d}{\mathrm dt}\varphi(r+ct,t)\,\mathrm dt
 \right)\mathrm dr=0.
 \end{split}
\]
测试函数具有紧支集；换元后 \(r\) 与 \(t\) 也限于一个有界集内，故 Fubini 换序合法。内层积分为零，因为测试函数在 \(t=0\) 附近及充分大的 \(t\) 处都消失。

还需验证弱导数确实是函数，才能使用本章的强解定义。对固定 \(t\)，在 \(x=ct\) 两侧分部积分，连续性使函数值的边界项相消，得到
\[
 \int_Q u\varphi_x\,\mathrm dx\,\mathrm dt
 =-\int_Q\operatorname{sgn}(x-ct)\varphi\,\mathrm dx\,\mathrm dt.
\]
对固定 \(x\)，在可能落入 \((0,\infty)\) 的折点 \(t=x/c\) 两侧作同样计算，得到
\[
 \int_Q u\varphi_t\,\mathrm dx\,\mathrm dt
 =c\int_Q\operatorname{sgn}(x-ct)\varphi\,\mathrm dx\,\mathrm dt.
\]
因此 \(u_x=\operatorname{sgn}(x-ct)\)、\(u_t=-c\operatorname{sgn}(x-ct)\) 是一阶弱导数；\(\operatorname{sgn}(0)\) 的取值不影响等价类。函数 \(u\) 局部有界，这两个导数也有界，故 \(u\in W^{1,\infty}_{\mathrm{loc}}(Q)\)。特别地，它属于每个 \(W^{1,p}_{\mathrm{loc}}(Q)\)，其中 \(1\leq p\leq\infty\)，包括 \(H^1_{\mathrm{loc}}(Q)\)。用上述弱导数代入可得 \(u_t+cu_x=0\) 几乎处处，所以它也是本章所定义的一阶强解。这与折线上的经典导数不存在并不矛盾：强解使用的是弱导数和几乎处处等式。

最后，三角不等式给出
\[
 \bigl||x-ct|-|x|\bigr|\leq|c|t.
\]
上界与 \(x\) 无关，因此当 \(t\downarrow0\) 时，\(u(\cdot,t)\to|\cdot|\) 在每个紧区间上一致成立。
\end{solution}

% 自拟综合题：在非平面曲面上同时计算面积、旋度通量与边界环流，检查诱导方向。
% 先为题头及首个公式预留空间，使超链接锚点与实际题头位于同页。
\Needspace{6\baselineskip}
\begin{optionalexercise}\label{b2:ex:01-paraboloid-stokes}
令
\[
 S=\{\,(x,y,z)\in\mathbb R^3\mid z=1-x^2-y^2,\ x^2+y^2\leq1\,\},
 \qquad F(x,y,z)=\left(-\frac y2,\frac x2,xy\right).
\]
曲面 \(S\) 取向上的单位法向 \(n\)，边界取它诱导的单位切向 \(\tau\)。
\begin{enumerate}
 \item\label{b2:item:01-paraboloid-area} 用参数图 \(X(s,t)=(s,t,1-s^2-t^2)\)、\(s^2+t^2\leq1\)，求 \(n\)、面积元 \(\mathrm dS\) 和 \(S\) 的面积。
 \item\label{b2:item:01-paraboloid-curl} 直接计算旋度通量 \(\int_S\operatorname{curl}F\cdot n\,\mathrm dS\)，并指出计算中单位法向的分母如何与面积元相消。
 \item\label{b2:item:01-paraboloid-boundary} 确定 \(\partial S\) 的诱导方向，独立计算 \(\int_{\partial S}F\cdot\tau\,\mathrm d\ell\)，与上一问比较。
 \item\label{b2:item:01-paraboloid-reversal} 反转曲面法向时，诱导边界方向与上述两个积分分别如何改变？如果只反转其中一种方向，原来的 Stokes 等式是否仍成立？
\end{enumerate}
\end{optionalexercise}
\begin{solution}
参数图在闭圆盘邻域光滑，且
\[
 X_s\times X_t=(2s,2t,1).
\]
第三分量恒正，故这确是向上的定向，且参数图处处秩为 \(2\)。于是
\[
 n\bigl(X(s,t)\bigr)=\frac{(2s,2t,1)}{\sqrt{1+4s^2+4t^2}},
 \qquad \mathrm dS=\sqrt{1+4s^2+4t^2}\,\mathrm ds\,\mathrm dt.
\]
在参数平面用极坐标，得到
\[
 \mathcal H^2(S)=2\pi\int_0^1 r\sqrt{1+4r^2}\,\mathrm dr
 =\frac\pi6(5\sqrt5-1).
\]
这与投影圆盘面积 \(\pi\) 不同：曲面的斜率增加了面积元。

向量场 \(F\) 是全空间上的光滑场，直接计算得
\[
 \operatorname{curl}F=(x,-y,1).
\]
因此单位法向的分母与面积元中的同一平方根相消，旋度通量为
\[
 \begin{aligned}
 \int_S\operatorname{curl}F\cdot n\,\mathrm dS
 &=\int_{s^2+t^2\leq1}(2s^2-2t^2+1)\,\mathrm ds\,\mathrm dt\\
 &=\pi.
 \end{aligned}
\]
末一步使用圆盘在交换 \(s,t\) 下不变，故 \(s^2\)、\(t^2\) 的积分相等。

参数圆盘的正向边界在从上方看时为逆时针，正向图 \(X\) 把它送到曲面边界；因而可取
\[
 \gamma(\theta)=(\cos\theta,\sin\theta,0),\quad0\leq\theta\leq2\pi.
\]
它的速度长度为 \(1\)，且
\[
 F(\gamma(\theta))\cdot\gamma'(\theta)
 =\frac12\sin^2\theta+\frac12\cos^2\theta=\frac12.
\]
所以边界环流为 \(\int_0^{2\pi}\frac12\,\mathrm d\theta=\pi\)，与旋度通量一致。这里边界的第三坐标恒为零，故 \(F\) 的第三分量不参与环流，尽管它参与了曲面上的旋度计算。

把 \(n\) 改成 \(-n\) 时，外余法向 \(\eta\) 不变，诱导切向由 \(\tau=n\times\eta\) 变成 \(-\tau\)。因此两侧积分同时由 \(\pi\) 变成 \(-\pi\)，Stokes 公式仍成立。若只反转法向或只反转边界方向，则只有一侧变号，所得 \(-\pi\) 与 \(\pi\) 不相等；这不是 Stokes 公式失效，而是两种方向已经不相容。
\end{solution}
